% docs/translation/latex/chapters/06_catalogue.tex - chapter of the LaTeX edition of the book "translation_report" (SIMC++ documentation).
%   Converted once from docs/translation/markdown/06_catalogue.md by tools/md2latex.py; from then on THIS FILE is the source of the
%   LaTeX edition and may be edited by hand. Re-running tools/md2latex.py --write-sources for this
%   chapter OVERWRITES the hand edits (see docs/README.md).
%   Compile: cmake --build <build> --target pdf, or tools/build_docs_pdf.py (figures, xelatex passes,
%   makeindex with docs/common/latex/simcpp.ist); by hand, see docs/README.md. XeLaTeX, class
%   docs/common/latex/simcpp.cls (packages: fontspec, tikz, xcolor, hyperref, listings, longtable, booktabs,
%   titlesec, titletoc, fancyhdr, makeidx, preview, ...); fonts TeX Gyre Pagella, TeX Gyre Heros, DejaVu Sans
%   Mono. The figures are built first from docs/{manual,translation}/figures/*.tex (preview.sty) into PDFs.
\chapter{\texorpdfstring{Catalogue of Fortran bugs, quirks and deviations (generated)}{Catalogue of Fortran bugs, quirks and deviations (generated)}}\label{translation_report-06u5fcatalogue}
\simcabstract{Generated by \code{docs/\allowbreak{}translation/\allowbreak{}tools/\allowbreak{}gen\_\allowbreak{}catalogue.\allowbreak{}py} from the Category tables of FORTRAN\_\allowbreak{}DIFFERENCES.md\footnote{From SIMC to SIMC++, ``Differences between SIMC++ and the Fortran SIMC'' (\code{docs/\allowbreak{}translation/\allowbreak{}markdown/\allowbreak{}FORTRAN\_\allowbreak{}DIFFERENCES.\allowbreak{}md})} (§2-§6: the single source that the code markers and the \code{fortran\_\allowbreak{}diff} tests cite) and cross-checked against the finding ids of \code{project\_\allowbreak{}control/\allowbreak{}DECISIONS.\allowbreak{}md} (§6.4). Do not edit by hand. The per-reaction measurements of the effect of the corrections (L2 in both modes) are in FORTRAN\_\allowbreak{}DIFFERENCES.md §7-§17.}
\section{\texorpdfstring{Categories}{Categories}}\label{translation_report-06u5fcatalogue:61-categories}
\begin{simctable}{small}{3}{>{\raggedright\arraybackslash}p{\dimexpr 0.105\linewidth-2\tabcolsep\relax}>{\raggedright\arraybackslash}p{\dimexpr 0.733\linewidth-2\tabcolsep\relax}>{\raggedright\arraybackslash}p{\dimexpr 0.094\linewidth-2\tabcolsep\relax}}
%
\toprule\rowcolor{simchead}\simcth{Category} & \simcth{meaning} & \simcth{entries} \\\midrule\endfirsthead
\multicolumn{3}{@{}l}{\simccontinued}\\\toprule\rowcolor{simchead}\simcth{Category} & \simcth{meaning} & \simcth{entries} \\\midrule\endhead
\midrule\multicolumn{3}{r@{}}{\simccontinues}\\\endfoot
\bottomrule\endlastfoot
A & clear Fortran bug: corrected by default, Fortran kept with -{}-fortran-compat & 71 \\
A-NP & clear bug in a path that was not ported (fix when porting) & 0 \\
B & ambiguous, maybe intended: Fortran kept in both modes (expert list) & 72 \\
C & open question with the user (acceptance mode): Fortran kept in both modes & 3 \\
D & faithful Fortran behaviour, not a bug: kept in both modes & 93 \\
DEV & deliberate design deviation of SIMC++ (not a Fortran bug): SIMC++ behaviour in both modes & 15 \\
\end{simctable}
Columns: \code{id} (decision and finding id, \code{D-{}xxx\ Fn}; the decision text is in DECISIONS.md and the finding in \code{reference/\allowbreak{}P*.\allowbreak{}F*\_\allowbreak{}report.\allowbreak{}md}), Fortran location (file:\allowbreak{}line),\allowbreak{} the Fortran behaviour, what SIMC++ does in the default (corrected) mode and with \code{-{}-{}fortran-{}compat}\simcidx{Command-line options!fortran-compat@\code{-{}-{}fortran-{}compat}}, the effect on outputs, and the test (Category A), the question (Category B), the marker (Category C) or the reason it is kept (Category D).

\section{\texorpdfstring{The catalogue}{The catalogue}}\label{translation_report-06u5fcatalogue:62-the-catalogue}
\subsection*{\texorpdfstring{Category A: clear Fortran bugs (corrected by default)}{Category A: clear Fortran bugs (corrected by default)}}\label{translation_report-06u5fcatalogue:category-a-clear-fortran-bugs-corrected-by-default}
\begin{landscape}
\begin{simctable}{scriptsize}{8}{l>{\raggedright\arraybackslash}p{\dimexpr 0.096\linewidth-2\tabcolsep\relax}>{\raggedright\arraybackslash}p{\dimexpr 0.114\linewidth-2\tabcolsep\relax}>{\raggedright\arraybackslash}p{\dimexpr 0.163\linewidth-2\tabcolsep\relax}>{\raggedright\arraybackslash}p{\dimexpr 0.122\linewidth-2\tabcolsep\relax}>{\raggedright\arraybackslash}p{\dimexpr 0.135\linewidth-2\tabcolsep\relax}>{\raggedright\arraybackslash}p{\dimexpr 0.171\linewidth-2\tabcolsep\relax}>{\raggedright\arraybackslash}p{\dimexpr 0.085\linewidth-2\tabcolsep\relax}}
\caption{Category A: clear Fortran bugs (corrected by default)}\\%
\toprule\rowcolor{simchead}\simcth{\#} & \simcth{id} & \simcth{Fortran location} & \simcth{Fortran behaviour} & \simcth{default (corrected) mode} & \simcth{\code{-{}-{}fortran-{}compat}} & \simcth{effect} & \simcth{test} \\\midrule\endfirsthead
\multicolumn{8}{@{}l}{\simccontinued}\\\toprule\rowcolor{simchead}\simcth{\#} & \simcth{id} & \simcth{Fortran location} & \simcth{Fortran behaviour} & \simcth{default (corrected) mode} & \simcth{\code{-{}-{}fortran-{}compat}} & \simcth{effect} & \simcth{test} \\\midrule\endhead
\midrule\multicolumn{8}{r@{}}{\simccontinues}\\\endfoot
\bottomrule\endlastfoot
1 & D-006 pass\_cuts & simc.f:237 & the hadron delta upper cut uses \code{SPedge\%\allowbreak{}e\%\allowbreak{}delta\%\allowbreak{}max}\simcidx{Parameters!SPedge.e.delta.max@\code{SPedge\%e\%delta\%max}} & \code{SPedge\%\allowbreak{}p\%\allowbreak{}delta\%\allowbreak{}max}\simcidx{Parameters!SPedge.p.delta.max@\code{SPedge\%p\%delta\%max}} (\code{driver::\allowbreak{}PassCuts}) & Fortran (original behaviour) & \code{pass\_\allowbreak{}cuts} → npasscuts, \code{wtcontribute} (.hist wtcontr), resolution sums; with \code{hard\_\allowbreak{}cuts}\simcidx{Parameters!hard_cuts@\code{hard\_cuts}} also the accepted events. Only events with recon hadron delta between the two limits are affected. K3 (HMS hadron at 423 MeV/c, limits 10 \% vs 22 \%): 42 of 200000 accepted events (norad) and 722 of 200000 (rad) no longer pass, .hist wtcontr −0.013 \% / −0.230 \% and the RECON SUMMARY resolutions change in the 4th digit (the Fortran trees give the same fractions: 27 and 755 events). P1.2, K2, K4: 0 events (no hadron recon delta in the window; in K2 the hadron limit is the larger one). & \code{fortran\_\allowbreak{}diff\_\allowbreak{}D006\_\allowbreak{}pass\_\allowbreak{}cuts} \\
2 & D-006 kaon decay, fixed upstream in 1fae2d7 (D-081) & shared/\allowbreak{}transp.\allowbreak{}f:\allowbreak{}158 & K → ππ decay in the first half of an element gives the product the kaon mass (\code{m\_\allowbreak{}final\ =\ Mk}) & \code{Mpi}, as transp.\allowbreak{}f:\allowbreak{}251 and project.\allowbreak{}f:\allowbreak{}88. Since U.R \code{-{}-{}fortran-{}compat}\simcidx{Command-line options!fortran-compat@\code{-{}-{}fortran-{}compat}} follows upstream 59ccdf6 (Mpi too); the 294f662 \code{Mk} only with \code{-{}-{}fortran-{}compat\ -{}-{}compat-{}target\ 294f662}\simcidx{Command-line options!fortran-compat@\code{-{}-{}fortran-{}compat}} & \textbf{fixed upstream} in 1fae2d7: as the default (upstream 59ccdf6); the 294f662 Fortran with \code{-{}-{}compat-{}target\ 294f662}\simcidx{Command-line options!compat-target@\code{-{}-{}compat-{}target}} & kaons with \code{doing\_\allowbreak{}decay}\simcidx{Parameters!doing_decay@\code{doing\_decay}} only (P10): mass and momentum of the decay product. No effect on H(e,e'p). & \code{fortran\_\allowbreak{}diff\_\allowbreak{}D006\_\allowbreak{}kaon\_\allowbreak{}decay\_\allowbreak{}mass} \\
3 & D-011 .geni Q3 & simc.\allowbreak{}f:\allowbreak{}519-520 & the SHMS "Q3 … mid" column prints \code{shmsSTOP\_\allowbreak{}q2\_\allowbreak{}mid} & prints \code{shmsSTOP\_\allowbreak{}q3\_\allowbreak{}mid} & Fortran (original behaviour) & \code{.\allowbreak{}geni} text only (the summary JSON keeps \code{q2\_\allowbreak{}mid\_\allowbreak{}dup} for the comparison with Fortran .geni files). 0 = 0 in the P1.2 runs. & \code{fortran\_\allowbreak{}diff\_\allowbreak{}D011\_\allowbreak{}D018\_\allowbreak{}report\_\allowbreak{}formats} \\
4 & D-011/\allowbreak{}D-060/\allowbreak{}D-061 F8 .geni SOS slit & simc.f:477 & the SOS "Slit hor/\allowbreak{}vert/\allowbreak{}corners" line prints \code{s\allowbreak{}STOP\_\allowbreak{}slit\_\allowbreak{}vert,\ s\allowbreak{}STOP\_\allowbreak{}slit\_\allowbreak{}hor,\ s\allowbreak{}STOP\_\allowbreak{}slit\_\allowbreak{}oct} (horizontal and vertical swapped against the label and the HMS/SHMS lines) & prints slit\_hor, slit\_vert, slit\_oct (\code{FormatGeni}, P14.1) & Fortran (original behaviour) & \code{.\allowbreak{}geni} text of SOS runs only (eep\_\allowbreak{}h\_\allowbreak{}test\_\allowbreak{}pctnames: "1313 18 21" → "18 1313 21") & \code{l1\_\allowbreak{}simcrun\_\allowbreak{}ranlux\_\allowbreak{}sos\_\allowbreak{}eep\_\allowbreak{}h\_\allowbreak{}test\_\allowbreak{}pctnames} (both modes) \\
5 & D-061 F8 .geni HRS & simc.\allowbreak{}f:\allowbreak{}487-509 & the HRS-R/HRS-L "Slit hor/vert" lines print \code{slit\_\allowbreak{}vert,\ slit\_\allowbreak{}hor}; the HRS blocks have no "Successes" line and the HRS-L block no trailing blank line & slit\_hor, slit\_vert; a "Successes" line and a blank line after each HRS block (\code{FormatGeni}, P14.2). The TAB inside the label 'S1, S2, Cal' is kept & Fortran (original behaviour) & \code{.\allowbreak{}geni} text of HRS runs only & \code{l1\_\allowbreak{}simcrun\_\allowbreak{}ranlux\_\allowbreak{}hrs\_\allowbreak{}Le\allowbreak{}Rp\_\allowbreak{}rad\_\allowbreak{}cap} (both modes) \\
6 & D-061 F1 HRS transp lengths & mc\_hrsl.f / mc\_\allowbreak{}hrsr.\allowbreak{}f:\allowbreak{}248,\allowbreak{} 259, 307, 318, 388, 455, 466 & the transport lengths are written as \code{62.\allowbreak{}75333333d0} ...: under -fdefault-real-8 these are REAL(16) constants passed to the real*8 dummy \code{zd} of transp, which reads their first 8 bytes: zd = −1.34e−41, −1.25e−127, −3.13e+232, 3.02e+189, −1.46e−71, −1.01e−65, −1.24e+107 (classes 2, 3, 5, 6, 8, 10, 11). transp adds zd to the pathlength (FP path = −3.134e+232 for every ray past Q2), uses it in the decay test and prints 'PROBLEM WITH TRANSFORMATION' at every call & the intended lengths (\code{Mc\allowbreak{}Hrs::\allowbreak{}transp\allowbreak{}Length}\simcidx{Classes and interfaces!McHrs@\code{McHrs}}) & Fortran (original behaviour) & HRS FP path (ntuple), decay probabilities in Q1/Q2/D1/Q3 with doing\_\allowbreak{}decay (every ray reaching the hut has decayed), the HRS kaon and SIDIS-kaon survival weight (physics\_\allowbreak{}kaon.\allowbreak{}f:\allowbreak{}159-163,\allowbreak{} semi\_\allowbreak{}physics.\allowbreak{}f:\allowbreak{}621-626: +Inf for every event, reference/\allowbreak{}runs/\allowbreak{}p14f1\_\allowbreak{}kaon\_\allowbreak{}hrs\_\allowbreak{}probe),\allowbreak{} the stdout of every HRS run (\textasciitilde{}420 MB per 1e5 events) & \code{l1\_\allowbreak{}mc\_\allowbreak{}hrs*} (compat, bitwise), \code{hrs\_\allowbreak{}selftest}, \code{fortran\_\allowbreak{}diff\_\allowbreak{}D061\_\allowbreak{}F1\_\allowbreak{}hrs\_\allowbreak{}transp\_\allowbreak{}lengths} (both modes, kaon weight) \\
7 & D-061 F1e HRS VDC rotation & mc\_\allowbreak{}hrsl\_\allowbreak{}hut.\allowbreak{}f / mc\_\allowbreak{}hrsr\_\allowbreak{}hut.\allowbreak{}f:\allowbreak{}293,\allowbreak{} 356 & \code{call\ rotate\_\allowbreak{}haxis(\allowbreak{}45.\allowbreak{}0d0,\allowbreak{}xt,\allowbreak{}yt)}: 45.0d0 is REAL(16), the real\emph{8 dummy reads +0.0 (low 8 bytes), so the VDC acceptance test (|x| ≤ 105.6, |y| ≤ 14.4) is done WITHOUT the 45° rotation (measured in reference/\allowbreak{}functions/\allowbreak{}hrs}/hut\_hrs*: xt = xs, yt = ys) & rotates by 45° as written (the code comment asks 'CHECK SIGN AND SIZE OF ROTATION') & Fortran (original behaviour) & HRS dc1/dc2 stops and the accepted HRS phase space & \code{l1\_\allowbreak{}hut\_\allowbreak{}hrs*} (compat, bitwise), \code{hrs\_\allowbreak{}selftest}, \code{fortran\_\allowbreak{}diff\_\allowbreak{}D061\_\allowbreak{}F1e\_\allowbreak{}hrs\_\allowbreak{}vdc\_\allowbreak{}rotation} \\
8 & D-062 G1 π0 photon decay & simc.\allowbreak{}f:\allowbreak{}1518-1522,\allowbreak{} 1535-1539 & the π0 photons go to mc\_calo with m2 = Mh2 = Mpi0² and decay\_flag = doing\_\allowbreak{}decay: with doing\_\allowbreak{}decay = 1 the first "photon decay" STOPs in project.\allowbreak{}f:\allowbreak{}76-90 ('cannot decay particle with mass=  134.98'; reference/\allowbreak{}runs/\allowbreak{}p14f2\_\allowbreak{}excl\_\allowbreak{}h\_\allowbreak{}dec1\_\allowbreak{}probe) & photons are massless and stable (m2 = 0, decay\_flag false; \code{driver::\allowbreak{}Pizero\allowbreak{}Calorimeter}\simcidx{Classes and interfaces!PizeroCalorimeter@\code{PizeroCalorimeter}}); the corrected set-up accepts doing\_\allowbreak{}decay with doing\_\allowbreak{}pizero (no effect) & Fortran (original behaviour) & every π0 run with doing\_\allowbreak{}decay = 1 (Fortran: no output at all) & \code{fortran\_\allowbreak{}diff\_\allowbreak{}D062\_\allowbreak{}G1\_\allowbreak{}pi0\_\allowbreak{}photon\_\allowbreak{}decay}, \code{calo\_\allowbreak{}selftest} \\
9 & D-062 G2 calo electron decay & simc.\allowbreak{}f:\allowbreak{}1761-1765 & the calorimeter as electron arm passes decay\_flag = doing\_\allowbreak{}decay (every other electron arm .false.): an electron "decay" STOPs (mass 0.51) & decay\_flag false (\code{MonteCarlo}\simcidx{Classes and interfaces!MonteCarlo@\code{MonteCarlo}}) & Fortran (original behaviour) & electron arm 7/8 with doing\_\allowbreak{}decay = 1 & \code{fortran\_\allowbreak{}diff\_\allowbreak{}D062\_\allowbreak{}G2\_\allowbreak{}calo\_\allowbreak{}electron\_\allowbreak{}decay} \\
10 & D-062 G3 π0 path length & simc.\allowbreak{}f:\allowbreak{}1462,\allowbreak{} 1517-1544 & pathlen is set to 0 once; the second photon starts from the first photon's path length, so FP\%p\%path ≈ 2 × drift\_\allowbreak{}to\_\allowbreak{}cal (600 cm for 300) & each photon starts at 0; FP\%p\%path = the path of the photon reported in FP\%p & Fortran (original behaviour) & π0 ntuple: no FP path column in the pion/SIDIS layouts; the summary/\allowbreak{}kaon-like survival weights do not use it & \code{fortran\_\allowbreak{}diff\_\allowbreak{}D062\_\allowbreak{}G3\_\allowbreak{}pi0\_\allowbreak{}path\_\allowbreak{}length}, \code{calo\_\allowbreak{}selftest} \\
11 & D-067 T1 field opposite side & simc.f:127, 141 & \code{degrad*abs(\allowbreak{}targ\_\allowbreak{}Bphi-{}phi)-{}180.\allowbreak{}0\ .\allowbreak{}lt.\ .\allowbreak{}01} has no outer abs: every Δφ < 180.01° that is not "same side" is taken as "opposite side" (targ\_Bphi = 0, the default: p arm Δφ = 90° gives 92.7°, the e arm hits the error branch and keeps 0) & \code{abs(\allowbreak{}degrad*abs(\allowbreak{}Δφ)-{}180)\ <\ .\allowbreak{}01} (driver::\allowbreak{}TargetFieldArmAngles) & Fortran (original behaviour) & using\_\allowbreak{}tgt\_\allowbreak{}field runs with targ\_Bphi ≠ the arm φ ± 0/180° only (probe p15f2\_\allowbreak{}poltar\_\allowbreak{}bphi0: both arms error branch, field angles 0) & \code{fortran\_\allowbreak{}diff\_\allowbreak{}D067\_\allowbreak{}T1\_\allowbreak{}field\_\allowbreak{}opposite\_\allowbreak{}side} \\
12 & D-067 T2 field p-arm sign & simc.\allowbreak{}f:\allowbreak{}136-140 & p arm, same side, targ\_\allowbreak{}Bangle < θp: \code{+1*sin(\allowbreak{}phi)*(\allowbreak{}targ\_\allowbreak{}Bangle-{}theta)} = −sinφ·|B−θ| (the sign of the B ≥ θ branch), unlike the continuous e-arm formula (simc.\allowbreak{}f:\allowbreak{}122-126) & the e-arm formula +sinφ·(θ−B) & Fortran (original behaviour) & field on the hadron-arm side with targ\_\allowbreak{}Bangle < θp only (probe p15f2\_\allowbreak{}poltar\_\allowbreak{}bangle5: −7.7° → +7.7°) & \code{fortran\_\allowbreak{}diff\_\allowbreak{}D067\_\allowbreak{}T2\_\allowbreak{}field\_\allowbreak{}parm\_\allowbreak{}sign} \\
13 & D-067 T14 field + π0/rho columns & results\_\allowbreak{}write.\allowbreak{}f:\allowbreak{}154-185,\allowbreak{} 212-240 & with using\_\allowbreak{}tgt\_\allowbreak{}field the π0 photons are written to ntu(54..65) (semi 57..68), over the 8 polarized-angle columns, and the last 8 photon tags stay 0; rho + field never writes \code{mmnuc} (tag 67, always 0) & each value in its own tag (OUTPUT\_\allowbreak{}FORMAT.md §1g) & Fortran (original behaviour) & ntuple only, π0/rho runs with the field (probe p15f2\_\allowbreak{}poltar\_\allowbreak{}pi0: 73 columns) & \code{fortran\_\allowbreak{}diff\_\allowbreak{}D067\_\allowbreak{}T14\_\allowbreak{}field\_\allowbreak{}ntuple\_\allowbreak{}columns} \\
14 & D-011 .geni DC & simc.f:526 & label "DC1, DC2, Scin, Cal" with two values & label "DC1, DC2" & Fortran (original behaviour) & \code{.\allowbreak{}geni} text only (every SHMS run). & same \\
15 & D-018 offsets & simc.\allowbreak{}f:\allowbreak{}815-818 & \code{xptar/\allowbreak{}yptar\ offset} printed in rad under the unit \code{mr} & printed in mr (×1000) & Fortran (original behaviour) & \code{.\allowbreak{}hist} text, only when a spectrometer angle offset is non-zero (0 in all validation runs). & same \\
16 & D-018 L format & simc.\allowbreak{}f:\allowbreak{}897-898,\allowbreak{} 907, 927 & the integers \code{deForest\_\allowbreak{}flag}\simcidx{Parameters!deForest_flag@\code{deForest\_flag}}, \code{use\_\allowbreak{}first\_\allowbreak{}cer}\simcidx{Parameters!use_first_cer@\code{use\_first\_cer}}, \code{use\_\allowbreak{}expon}\simcidx{Parameters!use_expon@\code{use\_expon}} written with the L edit descriptor (T/F; −1 and +1 both "T") & written as integers (i2, i2, i3) & Fortran (original behaviour) & \code{.\allowbreak{}hist} text of every run (2 or 3 lines). D-018 quoted l.896/\allowbreak{}906/\allowbreak{}925; the lines are 898/\allowbreak{}907/\allowbreak{}927. & same \\
17 & D-018 delta title & simc.\allowbreak{}f:\allowbreak{}727-732 & the (e,e'p)π reaction title goes to unit 6, not to the report unit & written to the \code{.\allowbreak{}hist} & Fortran (original behaviour) & \code{.\allowbreak{}hist} text of Delta runs (P11). & same \\
18 & D-014 deForest & physics\_\allowbreak{}proton.\allowbreak{}f:\allowbreak{}84-86 & the \code{cos\_\allowbreak{}phi} warning is tested after \code{cos\_\allowbreak{}phi} was clamped to ±1, so it never prints & tested (and printed) before the clamp & Fortran (original behaviour) & warning text only; the cross section is unchanged. D/A(e,e'p): printed for the central event of the A(e,e'p) inputs of P8.F1 (cos\_phi = 1 + 3.5e-8). & \code{fortran\_\allowbreak{}diff\_\allowbreak{}D014\_\allowbreak{}deforest\_\allowbreak{}cos\_\allowbreak{}phi\_\allowbreak{}warning} \\
19 & D-015 F9 raddeg & shared/\allowbreak{}rotate\_\allowbreak{}haxis.\allowbreak{}f:\allowbreak{}41,\allowbreak{} rotate\_\allowbreak{}vaxis.\allowbreak{}f:\allowbreak{}40 & \code{raddeg\ =\ 0.\allowbreak{}017453292} (π/180 to 8 digits, rel. error 1.1e-8) & π/180 & Fortran (original behaviour) & the rotated aperture coordinates of the HMS dipole (±6°) and the SHMS apertures (±1.5°, …) move by \textasciitilde{}1e-8 relative; the transported track is not changed. 0 aperture decisions changed in the 10 validation runs (trees bitwise identical, §7). & \code{fortran\_\allowbreak{}diff\_\allowbreak{}D015\_\allowbreak{}rotate\_\allowbreak{}raddeg} \\
20 & D-017 extreme trip & target.\allowbreak{}f:\allowbreak{}501 & hadron section, pudding can / cryo17: \code{zz} from the ELECTRON angle \code{the\%\allowbreak{}min} & from \code{thp\%\allowbreak{}min} & Fortran (original behaviour) & \code{targ\%\allowbreak{}Eloss(\allowbreak{}3)\%\allowbreak{}max}, \code{teff(\allowbreak{}3)\%\allowbreak{}max} (→ hadron edge limits, .hist "Eloss/teff lo-hi"). Only for \code{targ\%\allowbreak{}can}\simcidx{Parameters!targ.can@\code{targ\%can}} 2/3 AND an angle > 45° (otherwise z is clamped to −L/2 either way). Example: can 2, hadron at 52–60°: 1.3565 → 1.3591 MeV. All validation runs use the beer can (can 1): no effect. & \code{fortran\_\allowbreak{}diff\_\allowbreak{}D017\_\allowbreak{}extreme\_\allowbreak{}trip\_\allowbreak{}hadron\_\allowbreak{}angle} \\
21 & D-021 coll slits & hms/\allowbreak{}mc\_\allowbreak{}hms\_\allowbreak{}coll.\allowbreak{}f:\allowbreak{}83-104 & \code{h\allowbreak{}STOP\_\allowbreak{}slit\_\allowbreak{}hor/\allowbreak{}vert/\allowbreak{}oct} incremented for each of the 20 steps outside the opening, also for rays that pass & at most once per ray and counter (the Fortran's unused \code{slit\_\allowbreak{}*\_\allowbreak{}flag} flags) & Fortran (original behaviour) & \code{.\allowbreak{}geni} HMS slit counters with \code{using\_\allowbreak{}HMScoll}\simcidx{Parameters!using_HMScoll@\code{using\_HMScoll}} and pions/muons (P9). Example (2000 pions): 9746/\allowbreak{}3955/\allowbreak{}10944 → 637/\allowbreak{}263/\allowbreak{}713. Transport, RNG and acceptance unchanged. & \code{fortran\_\allowbreak{}diff\_\allowbreak{}D021\_\allowbreak{}hms\_\allowbreak{}coll\_\allowbreak{}slit\_\allowbreak{}counters} \\
22 & D-028 F2 NaN & event.\allowbreak{}f:\allowbreak{}1641 & \code{dy\ =\ sqrt(\allowbreak{}1/\allowbreak{}cos\_\allowbreak{}dtheta**2-{}1.\allowbreak{}-{}dx**2)} of a rounding-negative argument → NaN & argument clamped to ≥ 0 (dy = 0) & Fortran (original behaviour) & H elastic: the hadron yptar of a try can be NaN; such an event carries NaN in every hadron-arm recon/fp value but a finite weight (Fortran K2 norad: 1 of 200000 events). With seed 4242: 1 failed try in P1.2 rad and 1 in K3 rad (\code{.\allowbreak{}gen.\allowbreak{}root} geni\_\allowbreak{}p\_\allowbreak{}yptar +1 entry in the bin at 0), 0 accepted events in all 10 runs (the C++ K2 norad has no NaN event with this seed). & \code{fortran\_\allowbreak{}diff\_\allowbreak{}D028\_\allowbreak{}spectrometer\_\allowbreak{}angles\_\allowbreak{}nan} (110 NaN rows of the L1 table) \\
23 & D-033 F3 Benhar SF & sf\_\allowbreak{}lookup.\allowbreak{}f:\allowbreak{}93 (\code{sf\_\allowbreak{}lookup\_\allowbreak{}diff}) & \code{SFd\ =\ SF/\allowbreak{}4/\allowbreak{}3.\allowbreak{}1415926535/\allowbreak{}(\allowbreak{}Pm*Pm)/\allowbreak{}5.\allowbreak{}0/\allowbreak{}20.\allowbreak{}0}: the interpolated BIN probability divided by the event's own Pm² and by fixed 5 MeV × 20 MeV/c widths: +Inf/NaN at Pm = 0, a 1/Pm² divergence below the first node (Pm < 9.9-10 MeV/c), widths wrong for benharsf\_\allowbreak{}3mod/\allowbreak{}4 (dPm 19.733, dEm 1/5/10) & the differential SF of every node, sfval/(4π Pmval² dPm dEm) with the node's widths from the file, interpolated with the same bilinear scheme (\code{nuclear::\allowbreak{}Sf\allowbreak{}Table::\allowbreak{}lookup\allowbreak{}Diff}\simcidx{Classes and interfaces!SfTable@\code{SfTable}}) & Fortran (original behaviour) & A(e,e'p) with \code{use\_\allowbreak{}benhar\_\allowbreak{}sf}\simcidx{Parameters!use_benhar_sf@\code{use\_benhar\_sf}} only: main\%\allowbreak{}SF\_\allowbreak{}weight,\allowbreak{} hence Weight and wtcontr. §8 has the measured effect (eep\_\allowbreak{}fe56\_\allowbreak{}bh\_\allowbreak{}rad\_\allowbreak{}A). & \code{fortran\_\allowbreak{}diff\_\allowbreak{}D033\_\allowbreak{}F3\_\allowbreak{}benhar\_\allowbreak{}sf\_\allowbreak{}diff} \\
24 & D-033 F5 PmPer & event.f:926 & \code{vertex\%\allowbreak{}Pm\allowbreak{}Per\ =\ sqrt(\allowbreak{}max(\allowbreak{}0,\ vertex\%\allowbreak{}Pm**2\ -{}\ Pm\allowbreak{}Par**2\ -{}\ Pm\allowbreak{}Oop**2))} is computed before vertex\%Pm is assigned (l.940/\allowbreak{}943),\allowbreak{} i.e. with the previous try's Pm & recomputed with this event's Pm after l.940/943 (\code{Eep\allowbreak{}Reaction\allowbreak{}Base::\allowbreak{}complete\allowbreak{}Ev}\simcidx{Classes and interfaces!EepReactionBase@\code{EepReactionBase}}) & Fortran (original behaviour) & D/A(e,e'p): vertex\%\allowbreak{}PmPer only. It is not written to the ntuple (the tree has the recon PmPer) nor used by the weight: no effect on the outputs. & \code{fortran\_\allowbreak{}diff\_\allowbreak{}D033\_\allowbreak{}F5\_\allowbreak{}pmper\_\allowbreak{}stale\_\allowbreak{}pm} \\
25 & D-033 F6 sf\_lookup state & sf\_\allowbreak{}lookup.\allowbreak{}f:\allowbreak{}160-184 & for Em == Emval(numEm) (and NaN Em) no branch assigns Em1/\allowbreak{}Em2/\allowbreak{}sf1/\allowbreak{}sf2: the SAVEd values of the previous call are used (NaN on the first call of a process) & Em == Emval(numEm) uses the last interval; no state is kept between calls; NaN Em gives 0 & Fortran (original behaviour) & Benhar SF (A(e,e'p) \code{use\_\allowbreak{}benhar\_\allowbreak{}sf}\simcidx{Parameters!use_benhar_sf@\code{use\_benhar\_sf}}; A(e,e'π/K) generate\_\allowbreak{}em,\allowbreak{} not ported yet): only events with Em exactly on the last node (measure zero with continuous Em). & \code{fortran\_\allowbreak{}diff\_\allowbreak{}D033\_\allowbreak{}F6\_\allowbreak{}sf\_\allowbreak{}lookup\_\allowbreak{}stale} \\
26 & D-033 F7 heavy Mrec & event.f:944 & \code{vertex\%\allowbreak{}Mrec\ =\ sqrt(\allowbreak{}vertex\%\allowbreak{}Emiss**2-{}vertex\%\allowbreak{}Pmiss**2)} without a guard: NaN Mrec/\allowbreak{}Em/\allowbreak{}Trec and success = .true. (a NaN Em even passes the VERTEXedge tests, comparisons with NaN being false) & complete\_\allowbreak{}ev fails (success = false) when Emiss² − Pmiss² < 0 or NaN & Fortran (original behaviour) & A(e,e'p) at extreme edges only (He-4/C-12 table rows with ν → 0 or E\_p = Ein); 0 events in the L2 runs (§8). & \code{fortran\_\allowbreak{}diff\_\allowbreak{}D033\_\allowbreak{}F7\_\allowbreak{}heavy\_\allowbreak{}mrec\_\allowbreak{}guard} \\
27 & D-025/D-032 SHMS spectr & shms/\allowbreak{}mc\_\allowbreak{}shms.\allowbreak{}f:\allowbreak{}335-347 & the "LSA tune" STOP is tested on the first call only & tested on every call (both modes) & Fortran (original behaviour) & none in valid runs. & \code{l1\_\allowbreak{}shms\_\allowbreak{}selftest} ("spectr != 5 → FortranStop") \\
28 & D-032 arm code 6 & dbase.\allowbreak{}f:\allowbreak{}895-897/\allowbreak{}915-917 (STOP), but dbase.\allowbreak{}f:\allowbreak{}250/\allowbreak{}259,\allowbreak{} init.\allowbreak{}f:\allowbreak{}176/\allowbreak{}199,\allowbreak{} target.\allowbreak{}f:\allowbreak{}97/\allowbreak{}217,\allowbreak{} results\_\allowbreak{}write.\allowbreak{}f:\allowbreak{}90,\allowbreak{} simc.\allowbreak{}f:\allowbreak{}510-511/\allowbreak{}905-906 still treat 6 as SHMS & 6 = retired SHMS LSA tune, the run STOPs in dbase & rejected when the input is read (InpReader, JSON \code{FromJson}, \code{BuildRunSetup}), both modes: "electron\_\allowbreak{}arm = 6 was the retired SHMS LSA tune; use 5 = SHMS"; exit status 2. One spectrometer table (\code{core/\allowbreak{}Spectrometer\allowbreak{}Id.\allowbreak{}h}) is used everywhere; no other code knows 6. The shms plug-in key \code{"spectr"} was removed. & Fortran (original behaviour) & none for valid inputs. & \code{fortran\_\allowbreak{}diff\_\allowbreak{}D032\_\allowbreak{}spectrometer\_\allowbreak{}id}, \code{cli\_\allowbreak{}arm\_\allowbreak{}code\_\allowbreak{}6\_\allowbreak{}rejected} \\
29 & D-038 F10 sigcm2 & physics\_\allowbreak{}pion.\allowbreak{}f:\allowbreak{}133-154 & \code{ntup\%\allowbreak{}sigcm2} (the MAID cross section) is assigned only for W < 2 GeV; for W ≥ 2 GeV it keeps the value of an earlier event (or of the central event) & \code{ntup\%\allowbreak{}sigcm2\ =\ 0} for W ≥ 2 GeV (MAID not evaluated) in \code{pion::\allowbreak{}peepi} (\code{main.\allowbreak{}ntup.\allowbreak{}sigcm2}) & Fortran (original behaviour) & pion runs only (P9): \code{ntup\%\allowbreak{}sigcm2} is not a column of the pion ntuple (results\_\allowbreak{}write.\allowbreak{}f:\allowbreak{}131-186 writes it only for kaons), and the cross section (\code{ntup\%\allowbreak{}sigcm}, \code{main\%\allowbreak{}sigcc}) is the same in both modes. L1 peepi table: 3139 of 5047 rows have W ≥ 2 GeV, 1137 of them with a non-zero stale value. & \code{fortran\_\allowbreak{}diff\_\allowbreak{}D038\_\allowbreak{}F10\_\allowbreak{}sigcm2\_\allowbreak{}stale} \\
30 & D-038 F1 Δ++ STOP, fixed upstream in 20d46fc (D-081) & dbase.\allowbreak{}f:\allowbreak{}823-826 & \code{if(\allowbreak{}doing\_\allowbreak{}hydpi\ .\allowbreak{}and.\ which\_\allowbreak{}pion.\allowbreak{}eq.\allowbreak{}3)\ stop\ 'Can''t\ do\ Delta+\ production\ from\ neutron\ for\ pi0'} does not test \code{doing\_\allowbreak{}pizero}\simcidx{Parameters!doing_pizero@\code{doing\_pizero}}: charged H(e,\allowbreak{}e'π−)Δ++ (set up at dbase.\allowbreak{}f:\allowbreak{}361-364,\allowbreak{} weighted 0.55 at event.\allowbreak{}f:\allowbreak{}1479-1484; infiles/\allowbreak{}fpi12\_\allowbreak{}deltapp.\allowbreak{}inp) can never run & the STOP only with \code{doing\_\allowbreak{}pizero}\simcidx{Parameters!doing_pizero@\code{doing\_pizero}} (SetupBuilder); charged π−Δ++ on H runs. Since U.R \code{-{}-{}fortran-{}compat}\simcidx{Command-line options!fortran-compat@\code{-{}-{}fortran-{}compat}} follows upstream 59ccdf6 (runs too; full-run identity \code{l1\_\allowbreak{}simcrun\_\allowbreak{}ranlux\_\allowbreak{}up\_\allowbreak{}upf1\_\allowbreak{}pion\_\allowbreak{}h\_\allowbreak{}wp3\_\allowbreak{}A}); the 294f662 STOP only with \code{-{}-{}fortran-{}compat\ -{}-{}compat-{}target\ 294f662}\simcidx{Command-line options!fortran-compat@\code{-{}-{}fortran-{}compat}} & \textbf{fixed upstream} in 20d46fc: as the default (upstream 59ccdf6); the 294f662 Fortran with \code{-{}-{}compat-{}target\ 294f662}\simcidx{Command-line options!compat-target@\code{-{}-{}compat-{}target}} & H(e,\allowbreak{}e'π−)Δ++ inputs only (none of the validation runs). The π0Δ0-from-the-neutron case still STOPs in both modes. & \code{fortran\_\allowbreak{}diff\_\allowbreak{}D038\_\allowbreak{}F1\_\allowbreak{}delta\_\allowbreak{}pp\_\allowbreak{}stop} (incl. a 50-event corrected run) \\
31 & D-038 F2 π0 from n on H & dbase.\allowbreak{}f:\allowbreak{}141-142 & π− from hydrogen is stopped, π0 "from the neutron" (\code{which\_\allowbreak{}pion\ =\ 1}, \code{doing\_\allowbreak{}pizero}\simcidx{Parameters!doing_pizero@\code{doing\_pizero}}) is accepted on hydrogen and runs with Mtar\_\allowbreak{}struck = Mn on a proton target & rejected at start-up (\code{std::\allowbreak{}invalid\_\allowbreak{}argument}, exit status 2, "hydrogen has no neutron") & Fortran (original behaviour) & such inputs only (L1 cfg 15 of P9.F1). & \code{fortran\_\allowbreak{}diff\_\allowbreak{}D038\_\allowbreak{}F2\_\allowbreak{}pi0\_\allowbreak{}neutron\_\allowbreak{}on\_\allowbreak{}h} \\
32 & D-038 F5 CENTRAL.sigcc & simc.\allowbreak{}f:\allowbreak{}1236-1290 & calculate\_\allowbreak{}central calls only complete\_\allowbreak{}recon\_\allowbreak{}ev(vertex0): the central peepi sees main0\%\allowbreak{}epsilon = t = phi\_pq = 0 (σ\_L ∝ |t| = 0, ε = 0), so the \code{.\allowbreak{}hist} CENTRAL.sigcc is not the cross section at the central kinematics & main0 gets the central W, ε, θpq, t, tmin, Q², φpq (the complete\_\allowbreak{}ev formulas on vertex0; \code{IReaction::\allowbreak{}prepare\allowbreak{}Central}\simcidx{Classes and interfaces!IReaction@\code{IReaction}}, \code{PionProduction}\simcidx{Classes and interfaces!PionProduction@\code{PionProduction}}) & Fortran (original behaviour) & \code{.\allowbreak{}hist} CENTRAL.sigcc of pion runs only (not used by the weights or the normalisation). K1 π+ H: 0.839106E-06 → 0.925379E-05 (×11.03). §9 lists every L2 run. & \code{fortran\_\allowbreak{}diff\_\allowbreak{}D038\_\allowbreak{}F5\_\allowbreak{}central\_\allowbreak{}sigcc} \\
33 & D-038 F15 .hist which\_pion 11 & simc.\allowbreak{}f:\allowbreak{}750-755 & the final-state line is printed for which\_pion 0/1/10/2/3 but not for 11 (coherent π−) & \code{Final\ State\ is\ A\ +\ pi} also for 11 & Fortran (original behaviour) & \code{.\allowbreak{}hist} text of coherent π− runs. & \code{fortran\_\allowbreak{}diff\_\allowbreak{}D038\_\allowbreak{}F15\_\allowbreak{}hist\_\allowbreak{}which\_\allowbreak{}pion\_\allowbreak{}11} \\
34 & D-037 F1 SHMS coll slits & shms/\allowbreak{}mc\_\allowbreak{}shms\_\allowbreak{}coll.\allowbreak{}f:\allowbreak{}79,\allowbreak{} 86, 93 & \code{shms\allowbreak{}STOP\_\allowbreak{}slit\_\allowbreak{}hor/\allowbreak{}vert/\allowbreak{}oct} incremented for each of the 20 steps outside the opening (the SHMS copy of D-021) & at most once per ray and counter, as the HMS (P6.4; the P7.5 correction covered only mc\_\allowbreak{}hms\_\allowbreak{}coll) & Fortran (original behaviour) & \code{.\allowbreak{}geni} SHMS slit counters with \code{using\_\allowbreak{}SHMScoll}\simcidx{Parameters!using_SHMScoll@\code{using\_SHMScoll}} and pions/\allowbreak{}muons. Capture p6f2\_\allowbreak{}pi\_\allowbreak{}shms\_\allowbreak{}cap (3912 SHMS pions): 5646/\allowbreak{}20/\allowbreak{}5833 → 386/1/391. Transport, RNG and acceptance unchanged. & \code{fortran\_\allowbreak{}diff\_\allowbreak{}D037\_\allowbreak{}F1\_\allowbreak{}shms\_\allowbreak{}coll\_\allowbreak{}slit\_\allowbreak{}counters} \\
35 & D-037 F2 coll stops & simc.\allowbreak{}f:\allowbreak{}462-532 (hSTOP\_\allowbreak{}coll mc\_\allowbreak{}hms.\allowbreak{}f:\allowbreak{}218,\allowbreak{} sSTOP\_coll mc\_\allowbreak{}shms.\allowbreak{}f:\allowbreak{}518) & the collimator stops (absorbed pions) are written to no output: the \code{.\allowbreak{}geni} stop counts do not add up to the trials & \code{.\allowbreak{}geni} line \code{Collimator\ stops\ (\allowbreak{}pi/\allowbreak{}mu)} after the slit line of HMS/SHMS whenever \code{using\_\allowbreak{}HMScoll}\simcidx{Parameters!using_HMScoll@\code{using\_HMScoll}}/\code{using\_\allowbreak{}SHMScoll}\simcidx{Parameters!using_SHMScoll@\code{using\_SHMScoll}} is on (or the count is non-zero); the summary JSON has always had \code{coll} & Fortran (original behaviour) & \code{.\allowbreak{}geni} text only. Captures: 1253 of 3932 HMS pions (p6f2\_\allowbreak{}pi\_\allowbreak{}hms\_\allowbreak{}cap),\allowbreak{} 258 of 3912 SHMS pions (p6f2\_\allowbreak{}pi\_\allowbreak{}shms\_\allowbreak{}cap) become visible. The corrected slit counters (D-021) count rays that touched the collimator, not stops, so even then the lines add up only without them. & \code{fortran\_\allowbreak{}diff\_\allowbreak{}D037\_\allowbreak{}F2\_\allowbreak{}coll\_\allowbreak{}stops\_\allowbreak{}geni} \\
36 & D-037 F4 ranged out & hms/\allowbreak{}mc\_\allowbreak{}hms\_\allowbreak{}coll.\allowbreak{}f:\allowbreak{}121-124,\allowbreak{} shms/\allowbreak{}mc\_\allowbreak{}shms\_\allowbreak{}coll.\allowbreak{}f:\allowbreak{}123-126 (with enerloss\_\allowbreak{}new.\allowbreak{}f:\allowbreak{}75-77) & \code{if\ (\allowbreak{}epart.\allowbreak{}lt.\allowbreak{}sqrt(\allowbreak{}m2))\ goto\ 500} is unreachable: enerloss\_\allowbreak{}new clamps the loss to (E−m)−1e-7, so a particle that ranges out continues with p ≈ 4.6e-3 MeV/c (path up to 2e10 cm) & a clamped loss (\code{Simc\allowbreak{}Material::\allowbreak{}last\allowbreak{}Energy\allowbreak{}Loss\allowbreak{}Clamped}\simcidx{Classes and interfaces!SimcMaterial@\code{SimcMaterial}}) stops the particle in the collimator (hSTOP\_\allowbreak{}coll/\allowbreak{}sSTOP\_\allowbreak{}coll) & Fortran (original behaviour) & only particles that range out in the tungsten (slow muons from kaon decays). 3 of 1602 rows of mc\_\allowbreak{}hms\_\allowbreak{}decay (720, 891, 990: Q1\_in after 4.7e9-2.1e10 cm → coll); 0 calls in the two pion captures. & \code{fortran\_\allowbreak{}diff\_\allowbreak{}D037\_\allowbreak{}F4\_\allowbreak{}coll\_\allowbreak{}ranged\_\allowbreak{}out} \\
37 & D-037 F7 decay drift & shared/\allowbreak{}project.\allowbreak{}f:\allowbreak{}113 & \code{tmpdrift\ =\ z\_\allowbreak{}drift\ -{}\ z\_\allowbreak{}decay/\allowbreak{}sqrt(\allowbreak{}1+dxdzs**2+dydzs**2)} with the slopes AFTER the decay: the track ends at z\_drift + z\_\allowbreak{}decay(1/\allowbreak{}√(1+s\_\allowbreak{}old²) − 1/\allowbreak{}√(1+s\_\allowbreak{}new²)) & the z already covered uses the parent slopes; the track ends at z\_drift & Fortran (original behaviour) & decays in \code{project} (doing\_\allowbreak{}decay,\allowbreak{} pions/\allowbreak{}kaons): positions, pathlen and downstream apertures of the decay product. project\_\allowbreak{}decay table: max |Δx|,|Δy| 17 cm, |Δpath| 26 cm (synthetic edge rows), 0.016 cm / 0.28 cm in the real capture steps. Captures (replay of every recorded call): HMS 141 of 3932 pion calls change pathlen (max 0.28 cm), 43 accepted events change their recon (max 0.0014 \% δ, 2.4e-5 xptar, 0.0013 cm ytar); SHMS 104 of 3912 (max 0.13 cm), 28 accepted (max 0.010 \% δ); no acceptance decision changed. & \code{fortran\_\allowbreak{}diff\_\allowbreak{}D037\_\allowbreak{}F7\_\allowbreak{}project\_\allowbreak{}decay\_\allowbreak{}drift} \\
38 & D-037 F9 decay non-π/K & shared/\allowbreak{}project.\allowbreak{}f:\allowbreak{}87-91,\allowbreak{} transp.\allowbreak{}f:\allowbreak{}160-165,\allowbreak{} 253-258; dbase.\allowbreak{}f:\allowbreak{}1147 & the decay of a particle that is neither a pion nor a kaon executes a bare \code{stop} (exit 0, no output); doing\_\allowbreak{}decay is not tied to the reaction, so e.g. H(e,e'p) with doing\_\allowbreak{}decay = 1 stops at the first proton decay & the set-up rejects doing\_\allowbreak{}decay with a hadron that is not π/K (\code{std::\allowbreak{}invalid\_\allowbreak{}argument}, exit status 2, "… is neither a pion nor a kaon …"); project/\allowbreak{}transp never check a non-π/K particle for decay (no random number). Compat: \code{simc::\allowbreak{}FortranStop}\simcidx{Classes and interfaces!FortranStop@\code{FortranStop}}, exit 3 & Fortran (original behaviour) & configurations with doing\_\allowbreak{}decay and a proton/ρ/π0 hadron only; none of the reference inputs. & \code{fortran\_\allowbreak{}diff\_\allowbreak{}D037\_\allowbreak{}F9\_\allowbreak{}decay\_\allowbreak{}non\_\allowbreak{}pi\_\allowbreak{}k}, \code{cli\_\allowbreak{}decay\_\allowbreak{}non\_\allowbreak{}pik} \\
39 & D-043 F1 fint ABI & physics\_\allowbreak{}kaon.\allowbreak{}f:\allowbreak{}262,\allowbreak{} 375, 492 (\code{real*4\ fint}) vs cern/\allowbreak{}fint.\allowbreak{}f:\allowbreak{}10 (implicit REAL = REAL(8) under -fdefault-real-8,\allowbreak{} returned as binary128 with -ff2c) & the REAL*4 callers read the low 32 bits of the binary128 result: every fint value is +0.0, so eekeek/\allowbreak{}eekeeks give sigma\_eep = 0 (ntuple column \code{saghai} = 0 for every event) and phspwght (dead code) 10000 & the value fint computes, rounded to REAL*4 as the declaration intends (\code{kaon::\allowbreak{}Saghai\allowbreak{}Options::\allowbreak{}fint\_\allowbreak{}value}\simcidx{Classes and interfaces!SaghaiOptions@\code{SaghaiOptions}}) & Fortran (original behaviour) & ntuple column \code{saghai} (ntup\%\allowbreak{}sigcm1) of every kaon run; the weight uses sig\_\allowbreak{}factorized and is unchanged. L1: eekeek/\allowbreak{}eekeeks/\allowbreak{}phspwght/\allowbreak{}peek \code{\_\allowbreak{}fintabi} tables bitwise (1151 of 1244 eekeek rows non-zero) & \code{fortran\_\allowbreak{}diff\_\allowbreak{}D043\_\allowbreak{}kaon\_\allowbreak{}model}, \code{l1\_\allowbreak{}*\_\allowbreak{}fintabi} \\
40 & D-043 F2 sigma0 file & dbase.\allowbreak{}f:\allowbreak{}662-677 with saghai\_\allowbreak{}sigma0.dat & line 1 of the file is a truncated fragment, so every 4-line record is read one line off: zsrff1..zsiff3 = the 5-value kinematics line, zsrff4..zsiff6 = the true ff1..ff3, the true ff4..ff6 are lost & the fragment is skipped and each record is read from its own lines (\code{kaon::\allowbreak{}Saghai\allowbreak{}Tables::\allowbreak{}load(\allowbreak{}dir,\ true)}); the 18 stray records at the start of the file (s = 2.898 in the slots of s = 2.851, Q2 = 0) are kept where the file puts them and reported (Category B below) & Fortran (original behaviour) & \code{saghai} column of Sigma0/\allowbreak{}Sigma-/\allowbreak{}coherent runs only (masked by F1 in the Fortran). 22800 of 22800 corrected ff1..3 values equal the compat ff4..6 & \code{fortran\_\allowbreak{}diff\_\allowbreak{}D043\_\allowbreak{}kaon\_\allowbreak{}model} \\
41 & D-043 F3 s grid & physics\_\allowbreak{}kaon.\allowbreak{}f:\allowbreak{}288-294 & eekeek interpolates on s nodes 2.6 + 0.3(i-1): nodes 9/10 = 5.0/5.3 GeV², but the 9th and 10th blocks read from saghai\_\allowbreak{}proton.\allowbreak{}dat are s = 5.1/5.4 & nodes 9/10 = 5.1/5.4 (\code{Saghai\allowbreak{}Options::\allowbreak{}proton\_\allowbreak{}grid\_\allowbreak{}fixed}\simcidx{Classes and interfaces!SaghaiOptions@\code{SaghaiOptions}}) & Fortran (original behaviour) & \code{saghai} column for s > 4.7 GeV² (W > 2.17 GeV) of K+Lambda runs: 465 of 1244 eekeek table rows, none with s ≤ 4.7 & \code{fortran\_\allowbreak{}diff\_\allowbreak{}D043\_\allowbreak{}kaon\_\allowbreak{}model} \\
42 & D-043 F5 survival p & physics\_\allowbreak{}kaon.\allowbreak{}f:\allowbreak{}160-162 & the survival probability of a kaon whose decay is not simulated uses beta*gamma of the CENTRAL momentum spec\%p\%P; p\_kaon = spec\%\allowbreak{}p\%\allowbreak{}P(1+delta/\allowbreak{}100) is computed and unused (against the comment l.145-146) & beta*gamma of p\_kaon (\code{Peek\allowbreak{}Config::\allowbreak{}survival\_\allowbreak{}p\_\allowbreak{}kaon}\simcidx{Classes and interfaces!PeekConfig@\code{PeekConfig}}) & Fortran (original behaviour) & EVENT WEIGHTS (and the \code{decdist} column = survival probability) of decay-OFF kaon runs; §10 has the L2 effect & \code{fortran\_\allowbreak{}diff\_\allowbreak{}D043\_\allowbreak{}kaon\_\allowbreak{}model} (1407 of 1428 decay-OFF peek rows change) \\
43 & D-043 F6 zaero & physics\_\allowbreak{}kaon.\allowbreak{}f:\allowbreak{}149-159 & no zaero branch for hadron\_arm 5 (SHMS): the SAVEd value of an earlier call is used (0 in a SIMC run, −82.8/−183 after an arm 2/3/4 call); the path correction \code{zaero*(\allowbreak{}1+dx²+dy²)} lacks the square root & zaero = 0 for every arm without an aerogel distance, \code{zaero*sqrt(\allowbreak{}1+dx²+dy²)} (\code{Peek\allowbreak{}Config::\allowbreak{}survival\_\allowbreak{}path\_\allowbreak{}fixed}\simcidx{Classes and interfaces!PeekConfig@\code{PeekConfig}}) & Fortran (original behaviour) & none for HMS/SHMS runs (zaero = 0 in both modes); 60 order-dependent SHMS rows of the peek table & \code{fortran\_\allowbreak{}diff\_\allowbreak{}D043\_\allowbreak{}kaon\_\allowbreak{}model} \\
44 & D-048 F1 Delta D/A STOP & dbase.\allowbreak{}f:\allowbreak{}779-788 (D-018) & the (e,e'p)π announcement tests \code{doing\_\allowbreak{}deutpi} / \code{doing\_\allowbreak{}hepi} (pion flags, never set for doing\_\allowbreak{}delta): every D or A Delta run STOPs with \code{I\ don't\ have\ ANY\ idea\ what\ (\allowbreak{}e,\allowbreak{}e'p)pi\ we're\ doing!!!} (exit 0) & tests \code{doing\_\allowbreak{}deutdelta} / \code{doing\_\allowbreak{}hedelta}: D(e,e'p)π and A(e,e'p)π run (\code{SetupBuilder}) & Fortran (original behaviour) & D/A Delta inputs only. The D/A branches of event.f (generate l.336-371, complete\_\allowbreak{}ev l.646-656, 949-971) and init.f (l.349-358) then run for the first time: \textbf{no Fortran reference exists for them}; they are validated by self-consistency only (energy-momentum balance of every event, \code{fortran\_\allowbreak{}diff\_\allowbreak{}D048\_\allowbreak{}F1\_\allowbreak{}delta\_\allowbreak{}DA\_\allowbreak{}stop}: D, He-3, C-12). & \code{fortran\_\allowbreak{}diff\_\allowbreak{}D048\_\allowbreak{}F1\_\allowbreak{}delta\_\allowbreak{}DA\_\allowbreak{}stop} \\
45 & D-048 F1 Delta D/A data & dbase.\allowbreak{}f:\allowbreak{}563-565,\allowbreak{} 594 & the momentum distribution (deut.dat, he3.dat, ...) and the Benhar spectral function (generate\_\allowbreak{}em,\allowbreak{} \code{Emval(\allowbreak{}numEm)}) are read for \code{doing\_\allowbreak{}deutpi/\allowbreak{}hepi/\allowbreak{}deutkaon/\allowbreak{}hekaon} only, not for \code{doing\_\allowbreak{}deutdelta/\allowbreak{}hedelta}, although generate and limits\_\allowbreak{}init use them for the Delta (unreachable in the Fortran because of the STOP above) & the delta flags load the same files as the pion (corrected mode only) & Fortran (original behaviour) & D/A Delta runs only (they exist in corrected mode only) & same \\
46 & D-048 F2 Delta tail 1 & radc.\allowbreak{}f:\allowbreak{}288-293 & the incident-electron tail limits are set for doing\_\allowbreak{}heavy,\allowbreak{} hyd\_elast, deuterium and \code{doing\_\allowbreak{}pion\ .\allowbreak{}or.\ doing\_\allowbreak{}kaon\ .\allowbreak{}or.\ doing\_\allowbreak{}rho\ .\allowbreak{}or.\ doing\_\allowbreak{}semi}, not for doing\_\allowbreak{}delta: Egamma\_\allowbreak{}min(1) = Egamma\_\allowbreak{}max(1) = 0, basicrad returns weight 0, and EVERY try with ntail = 1 fails (frac(1) = 0.43-0.46 of the tries at the P11.F1 kinematics) & doing\_\allowbreak{}delta uses CASE 4 (Egamma\_\allowbreak{}min(1) = 0, Egamma\_\allowbreak{}max(1) = gen\%\allowbreak{}sumEgen\%\allowbreak{}max − vertex\%\allowbreak{}e\%\allowbreak{}E): init.\allowbreak{}f:\allowbreak{}455-463 sets gen\%\allowbreak{}sumEgen for doing\_\allowbreak{}delta with the same "generated ELECTRON energy" branch as for the pion (\code{Radc\allowbreak{}Run::\allowbreak{}delta\_\allowbreak{}tail1\_\allowbreak{}case4}\simcidx{Classes and interfaces!RadcRun@\code{RadcRun}}) & Fortran (original behaviour) & EVENT RATES of rad Delta runs: Ntried and wtcontr. §11 has the L2 effect. With tail 1 radiating, generate\_\allowbreak{}rad calls complete\_\allowbreak{}ev again (radc.\allowbreak{}f:\allowbreak{}322),\allowbreak{} which draws a new root choice (Category B below). & \code{fortran\_\allowbreak{}diff\_\allowbreak{}D048\_\allowbreak{}F2\_\allowbreak{}delta\_\allowbreak{}tail1\_\allowbreak{}egamma} \\
47 & D-048 F3 sig\_blok args & physics\_\allowbreak{}delta.\allowbreak{}f:\allowbreak{}99-100 & \code{sig\_\allowbreak{}blok(\allowbreak{}thetacm,\allowbreak{}phicm,\allowbreak{}main\%\allowbreak{}t/\allowbreak{}1.\allowbreak{}d6,\allowbreak{}vertex\%\allowbreak{}q2/\allowbreak{}1.\allowbreak{}d6,\allowbreak{}s/\allowbreak{}1.\allowbreak{}d6,\allowbreak{}.\allowbreak{}.\allowbreak{}.\allowbreak{})}: under -fdefault-real-8 \code{1.\allowbreak{}d6} is REAL(16), so binary128 temporaries are passed to the REAL*8 dummies and sig\_blok reads their low 8 bytes: ntup\%sigcm (column \code{sigcm}) is meaningless (0 in \textasciitilde{}65 \%, negative in \textasciitilde{}19 \%, NaN/Inf in \textasciitilde{}0.4 \% of the events) & the doubles t, Q², s [GeV²] are passed (\code{delta::\allowbreak{}Delta\allowbreak{}Config::\allowbreak{}sigblok\_\allowbreak{}binary128\_\allowbreak{}quirk\ =\ false}): \code{sigcm} = the Blok π fit at the proton's CM angles (still not a Δ model, see F4) & Fortran (original behaviour) & ntuple column \code{sigcm} of Delta runs only; the weight is unchanged (F4). peedelta table: 3069 of 3077 rows change, 14 non-finite compat values become finite & \code{fortran\_\allowbreak{}diff\_\allowbreak{}D048\_\allowbreak{}F3\_\allowbreak{}delta\_\allowbreak{}sigblok\_\allowbreak{}args}, \code{l1\_\allowbreak{}peedelta} \\
48 & D-048 F6 Delta warning & dbase.\allowbreak{}f:\allowbreak{}184-186 & quote typo: \code{'WARNING:\ Delta\ cross\ section\ model",} + continuation \code{"\ only\ set\ up\ for\ proton\ target!'} is ONE string, printed with \code{",}, 79 blanks and \code{"} & the sentence without the quote debris & Fortran (original behaviour) & log text of D/A Delta runs & \code{fortran\_\allowbreak{}diff\_\allowbreak{}D048\_\allowbreak{}F6\_\allowbreak{}delta\_\allowbreak{}warning\_\allowbreak{}text} \\
49 & D-048 F9 CENTRAL.sigcc & simc.\allowbreak{}f:\allowbreak{}1236-1290 (as D-038 F5) & the central peedelta sees main0\%\allowbreak{}epsilon = t = phi\_pq = 0 & main0 gets the central kinematics (\code{Pion\allowbreak{}Production::\allowbreak{}prepare\allowbreak{}Central}\simcidx{Classes and interfaces!PionProduction@\code{PionProduction}}, inherited by \code{DeltaProduction}\simcidx{Classes and interfaces!DeltaProduction@\code{DeltaProduction}}) & Fortran (original behaviour) & \code{.\allowbreak{}hist} CENTRAL.sigcc of Delta runs only: D1 rad 0.227051E+01 → 0.167557E+02 (×7.38) & \code{fortran\_\allowbreak{}diff\_\allowbreak{}D048\_\allowbreak{}F9\_\allowbreak{}central\_\allowbreak{}sigcc\_\allowbreak{}delta} \\
50 & D-051 F1 sum\_sq = 0 & semi\_\allowbreak{}physics.\allowbreak{}f:\allowbreak{}508-513 & dsigdz = (...)/\allowbreak{}sum\_\allowbreak{}sq/\allowbreak{}zhad without the \code{sum\_\allowbreak{}sq\ >\ 0} guard of the old code (l.378-400, commented out): when every PDF is 0 (x\_Bj > 1 clamped to 1, Q < Alambda) the cross section and the event weight are 0/0 = NaN & dsigdz = 0 as the old guard (\code{Peepi\allowbreak{}XConfig::\allowbreak{}sum\_\allowbreak{}sq\_\allowbreak{}guard}\simcidx{Classes and interfaces!PeepiXConfig@\code{PeepiXConfig}}) & Fortran (original behaviour) & weight NaN → 0 for those events (the Fortran ΣWeight/\allowbreak{}.hist sums of such runs are NaN-contaminated if one occurs); 34 of the 3338 P12.F1 peepix rows & \code{fortran\_\allowbreak{}diff\_\allowbreak{}D051\_\allowbreak{}sidis} (34 NaN rows → 0, no other row), \code{sidis\_\allowbreak{}sanity} \\
51 & D-051 F2 negative Fermi x\_Bj & semi\_\allowbreak{}physics.\allowbreak{}f:\allowbreak{}247-248 → Ctq5Pdf.\allowbreak{}f:\allowbreak{}80-83 & with do\_fermi, x\_Bj = Q²/2/(efer ν − pfer p̂·q) is negative for large pfer along q (deut.dat up to 1190 MeV/c): Ctq5Pdf STOPs the run (exit 0, no .hist); a D run with do\_fermi = 1 cannot reach 1e5 events (P12.F1: stop at try 27848) & x\_Bj ≤ 0 returns 0 before the PDF calls (same early return as the threshold, \code{Peepi\allowbreak{}XConfig::\allowbreak{}nonpositive\_\allowbreak{}xbj\_\allowbreak{}zero}\simcidx{Classes and interfaces!PeepiXConfig@\code{PeepiXConfig}}); compat throws \code{FortranStop}\simcidx{Classes and interfaces!FortranStop@\code{FortranStop}} & Fortran (original behaviour) & D/A runs with do\_fermi = 1 complete; the affected events get weight 0 & \code{sidis\_\allowbreak{}sanity} (compat STOP, corrected 0); P12.F1 stops.txt \\
52 & D-051 F4 π0 FF/mass & semi\_\allowbreak{}physics.\allowbreak{}f:\allowbreak{}490-495,\allowbreak{} 227 & the doing\_\allowbreak{}pizero block ("just get averge of pi+ and pi-") assigns \code{d1\ =\ (\allowbreak{}yf\ +\ yu)\ /\ 2.} twice and never \code{db}: the d̄ fragmentation function keeps \code{db\ =\ u1} of l.487 (the π+ fit yf for doing\_\allowbreak{}hplus,\allowbreak{} the π− fit yu otherwise); and mhad = Mpi (charged) at l.227 while Mpi\_gev (l.149) and dbase.f Mh use Mpi0 & db = (yf+yu)/2 like u1, d1, ub, and mhad = Mpi0 (\code{sidis::\allowbreak{}Peepi\allowbreak{}XConfig::\allowbreak{}pizero\_\allowbreak{}fixed}\simcidx{Classes and interfaces!PeepiXConfig@\code{PeepiXConfig}}) & Fortran (original behaviour) & the d̄ term of the π0 SIDIS cross section; π0 SIDIS is refused until the calorimeter arm (P14), so no run is affected yet & \code{fortran\_\allowbreak{}diff\_\allowbreak{}D051\_\allowbreak{}sidis} (16 π0 rows of the P12.F1 peepix table: all change, no other row does), \code{sidis\_\allowbreak{}sanity} \\
53 & D-051 F6 survival p & semi\_\allowbreak{}physics.\allowbreak{}f:\allowbreak{}622-625 (as D-043 F5) & the survival probability of a hadron whose decay is not simulated uses β·γ of the CENTRAL momentum spec\%p\%P; p\_kaon = spec\%\allowbreak{}p\%\allowbreak{}P(1+δ/\allowbreak{}100) is computed and unused (against the comment l.607-608) & β·γ of p\_kaon (\code{Peepi\allowbreak{}XConfig::\allowbreak{}survival\_\allowbreak{}p\_\allowbreak{}hadron}\simcidx{Classes and interfaces!PeepiXConfig@\code{PeepiXConfig}}) & Fortran (original behaviour) & \code{decdist} of every decay-OFF SIDIS run; the WEIGHT of decay-OFF K± runs (event.\allowbreak{}f:\allowbreak{}1560 applies it for doing\_\allowbreak{}semika only) & \code{fortran\_\allowbreak{}diff\_\allowbreak{}D051\_\allowbreak{}sidis} (3268 of the 3281 decay-OFF peepix rows change, no other row), \code{sidis\_\allowbreak{}sanity} \\
54 & D-051 F6 zaero & semi\_\allowbreak{}physics.\allowbreak{}f:\allowbreak{}611-621 (as D-043 F6) & no zaero branch for hadron\_arm 5 (SHMS): the SAVEd value of an earlier call is used (0 in a SIMC run); the path correction \code{zaero*(\allowbreak{}1+dx²+dy²)} lacks the square root (HMS: zaero = −331.491 cm, so the path is ≈ 0.4 cm too short for dx = 0.05) & zaero = 0 for every arm without an aerogel distance, \code{zaero*sqrt(\allowbreak{}1+dx²+dy²)} (\code{Peepi\allowbreak{}XConfig::\allowbreak{}survival\_\allowbreak{}path\_\allowbreak{}fixed}\simcidx{Classes and interfaces!PeepiXConfig@\code{PeepiXConfig}}) & Fortran (original behaviour) & decay-OFF SIDIS runs with the HMS as hadron arm (sqrt); none for SHMS runs (0 in both modes) & \code{fortran\_\allowbreak{}diff\_\allowbreak{}D051\_\allowbreak{}sidis}, \code{sidis\_\allowbreak{}sanity} \\
55 & D-051 F3 pfermi NaN, fixed upstream in e4dd175 (D-081; since U.1 also in fortran\_\allowbreak{}compat) & results\_\allowbreak{}write.\allowbreak{}f:\allowbreak{}208-209 & the SIDIS ntuple column \code{pfermi} = pfer/\allowbreak{}1000·|d|/\allowbreak{}d with d = p̂\_\allowbreak{}fermi·q̂,\allowbreak{} without the \code{if\ (\allowbreak{}dummy.\allowbreak{}eq.\allowbreak{}0)\ dummy=\allowbreak{}1.\allowbreak{}e-{}20} of the pion block (l.143): 0/0 = NaN for every event with pfer = 0 (every H and A > 2 event) & the pion-block guard: pfermi = 0 (\code{FillSemiNtuple}) & \textbf{fixed upstream} in e4dd175: as the default (upstream 59ccdf6) & ntuple column \code{pfermi} of H and A > 2 SIDIS runs (NaN → 0); nothing else & \code{fortran\_\allowbreak{}diff\_\allowbreak{}D051\_\allowbreak{}F3\_\allowbreak{}F9\_\allowbreak{}semi\_\allowbreak{}pfermi\_\allowbreak{}pmper} (300 H events: compat NaN, corrected 0) \\
56 & D-051 F9 PmPer & event.f:922 before 952-954 (as D-033 F5) & the semi complete\_\allowbreak{}ev computes vertex\%\allowbreak{}PmPer from the previous try's vertex\%Pm, assigning Pm = Pmiss only afterwards & PmPer from this event's Pm (\code{Sidis\allowbreak{}Production::\allowbreak{}complete\allowbreak{}Ev}\simcidx{Classes and interfaces!SidisProduction@\code{SidisProduction}}) & Fortran (original behaviour) & vertex PmPer only (not an ntuple column of the SIDIS layout) & \code{fortran\_\allowbreak{}diff\_\allowbreak{}D051\_\allowbreak{}F3\_\allowbreak{}F9\_\allowbreak{}semi\_\allowbreak{}pfermi\_\allowbreak{}pmper} (compat 0 / corrected 300 of 300 events use the current Pm) \\
57 & D-051 F12 π0 title & simc.\allowbreak{}f:\allowbreak{}679-700 & the .hist header has no π0 line: a π0 SIDIS run is titled π+ or π− by doing\_\allowbreak{}hplus (dbase.f announces (e,e'pi0)X) & title (e,e'pi0)X & Fortran (original behaviour) & .hist header (π0 SIDIS itself is refused until the calorimeter arm, P14) & \code{fortran\_\allowbreak{}diff\_\allowbreak{}D051\_\allowbreak{}F12\_\allowbreak{}F14\_\allowbreak{}semi\_\allowbreak{}hist\_\allowbreak{}titles} \\
58 & D-051 F14 A ≥ 3 kaon & simc.\allowbreak{}f:\allowbreak{}701-717 & the .hist header has no A ≥ 3 branch for kaon SIDIS: \code{STOP\ I\ don't\ have\ ANY\ idea\ what\ A(\allowbreak{}e,\allowbreak{}e'k)X\ we're\ doing!!!} after the whole event loop (exit 0, .hist truncated to 102 bytes; .gen/.geni written) & \textbf{run}: the A(e,e'k±)X title is written and the report completes. Choice: the run itself is fully defined (dbase.\allowbreak{}f:\allowbreak{}768-773 announces it, peepiX handles A > 2 with slacemcfit/\allowbreak{}Rhad\_\allowbreak{}global and fDSS, the generation has no A-specific kaon branch); only the title branch is missing, so refusing the input would remove a working configuration & Fortran (original behaviour) & C-12 K± SIDIS runs complete with a full .hist & \code{fortran\_\allowbreak{}diff\_\allowbreak{}D051\_\allowbreak{}F12\_\allowbreak{}F14\_\allowbreak{}semi\_\allowbreak{}hist\_\allowbreak{}titles} (compat STOP, corrected .hist + a 100-event run) \\
59 & D-043 F13 CENTRAL.sigcc & simc.\allowbreak{}f:\allowbreak{}1236-1290 (as D-038 F5) & the central peeK sees main0\%\allowbreak{}epsilon = t = theta\_pq = phi\_pq = 0 & main0 gets the central W, ε, θpq, t, tmin, Q², φpq (\code{Pion\allowbreak{}Production::\allowbreak{}prepare\allowbreak{}Central}\simcidx{Classes and interfaces!PionProduction@\code{PionProduction}}, inherited by \code{KaonProduction}\simcidx{Classes and interfaces!KaonProduction@\code{KaonProduction}}) & Fortran (original behaviour) & \code{.\allowbreak{}hist} CENTRAL.sigcc of kaon runs only: K1 K+Λ H 0.143536E-06 → 0.704013E-06 (×4.90); §10 lists every L2 run & \code{fortran\_\allowbreak{}diff\_\allowbreak{}D043\_\allowbreak{}F13\_\allowbreak{}central\_\allowbreak{}sigcc\_\allowbreak{}kaon} \\
60 & D-057 F6 rho efer & rho\_\allowbreak{}physics.\allowbreak{}f:\allowbreak{}88-93 & peerho recomputes the COMMON efer on shell, sqrt(pfer²+Mtar\_\allowbreak{}struck²),\allowbreak{} and writes it back to /pfermi/: the off-shell branch tests the PION flags doing\_\allowbreak{}deutpi/\allowbreak{}doing\_\allowbreak{}hepi,\allowbreak{} never set in a rho run (1001/1001 D capture calls: efer 936.047 → 938.272) & peerho uses the COMMON efer of the generation and leaves it unchanged (\code{rho::\allowbreak{}Peerho\allowbreak{}Config::\allowbreak{}efer\_\allowbreak{}from\_\allowbreak{}common}\simcidx{Classes and interfaces!PeerhoConfig@\code{PeerhoConfig}}); which formula is right for A > 1 stays on the expert list (D-057) & Fortran (original behaviour) & D/³He rho weights (s, flux, CM boost); none for H (pfer = 0, efer = Mp in both). D: ×0.828 on a pfer = 200 MeV/c event; in a run pfer = 0 (F4), so only efer = Md − Mn vs Mp & \code{fortran\_\allowbreak{}diff\_\allowbreak{}D057\_\allowbreak{}rho\_\allowbreak{}kernels}, \code{l1\_\allowbreak{}peerho} (compat) \\
61 & D-057 F3 rho Mh2\_final & generate\_\allowbreak{}rho.\allowbreak{}f:\allowbreak{}80 & \code{Mh2\_\allowbreak{}final\ =\ Mh} (a mass, not a mass squared), 5695/5695 calls & Mh2\_final = Mh2 (\code{rho::\allowbreak{}Generate\allowbreak{}Rho\allowbreak{}Config::\allowbreak{}mh2\_\allowbreak{}final\_\allowbreak{}fixed}\simcidx{Classes and interfaces!GenerateRhoConfig@\code{GenerateRhoConfig}}) & Fortran (original behaviour) & none observable (simc.f:174 resets it, rho\_decay sets Mpi2 on every success) & \code{fortran\_\allowbreak{}diff\_\allowbreak{}D057\_\allowbreak{}rho\_\allowbreak{}kernels} \\
62 & D-057 F1 rho stale hadron angle & event.\allowbreak{}f:\allowbreak{}280-283,\allowbreak{} 325, 1034; generate\_\allowbreak{}rho.\allowbreak{}f:\allowbreak{}123-124 & generate draws no hadron slopes for doing\_rho, so physics\_\allowbreak{}angles converts the previous try's rho lab theta/phi (stored in vertex\%\allowbreak{}p\%\allowbreak{}xptar/\allowbreak{}yptar) and trip\_\allowbreak{}thru\_\allowbreak{}target(3) uses that angle for the rho energy loss and teff(3) (θ used: mean 1.06 rad vs the rho's 0.42 rad; ΔEloss(3) −6…+2.4 MeV, teff(3) −24…+58 \%) & trip\_\allowbreak{}thru\_\allowbreak{}target(3) uses the rho's own lab angles (\code{RhoProduction}\simcidx{Classes and interfaces!RhoProduction@\code{RhoProduction}}, corrected mode) & Fortran (original behaviour) & Eloss(3)/\allowbreak{}teff(3) of every event: the decay pion's energy-loss correction and multiple scattering in montecarlo & \code{fortran\_\allowbreak{}diff\_\allowbreak{}D057\_\allowbreak{}F1\_\allowbreak{}F2\_\allowbreak{}rho} (3975 capture tries, teff(3) changes in all) \\
63 & D-057 F2 rho\_decay after failure & event.\allowbreak{}f:\allowbreak{}420-422 & rho\_decay is called for every try that reaches l.420, also after a failed generate\_\allowbreak{}rad: it decays the PREVIOUS try's orig (3.4 \% of the H tries), draws 3–41 uniforms and may set Mh = Mpi & rho\_decay only after a successful generation. No weight changes; the RANDOM STREAM differs from the Fortran's from the first such try on (those 1+2k uniforms are not drawn), and Mh/\allowbreak{}Mh2/\allowbreak{}ntup\%\allowbreak{}rhotheta/\allowbreak{}orig keep their values on those tries & Fortran (original behaviour) & the random stream of a corrected run (no Fortran-identical replay in that mode); no physics quantity & \code{fortran\_\allowbreak{}diff\_\allowbreak{}D057\_\allowbreak{}F1\_\allowbreak{}F2\_\allowbreak{}rho} (94 capture tries: compat draws the decay uniforms, corrected not; outcome identical) \\
64 & D-057 F9 rho yptar sign & rho\_\allowbreak{}decay.\allowbreak{}f:\allowbreak{}119-124 & for a beam-left arm (2/4/5) and a decay pion on the far side of the beam (pyf ≤ 0), th\_inp = +(α + θ\_spec): yptar gets the wrong sign, and orig\%\allowbreak{}p\%\allowbreak{}theta/\allowbreak{}phi follow & th\_inp = −(α + θ\_spec) (\code{rho::\allowbreak{}Rho\allowbreak{}Decay\allowbreak{}Config::\allowbreak{}yptar\_\allowbreak{}sign\_\allowbreak{}fixed}\simcidx{Classes and interfaces!RhoDecayConfig@\code{RhoDecayConfig}}) & Fortran (original behaviour) & such pions have & yptar \\
65 & D-057 F20 rho NaN & rho\_\allowbreak{}physics.\allowbreak{}f:\allowbreak{}388-389 & the 'crazy sigma' guard (> 1e10 or < 0 → 0) lets NaN through (s < 0 edge row) & NaN → 0 (\code{rho::\allowbreak{}Peerho\allowbreak{}Config::\allowbreak{}nan\_\allowbreak{}guard}\simcidx{Classes and interfaces!PeerhoConfig@\code{PeerhoConfig}}) & Fortran (original behaviour) & edge cases only (no NaN in the L2 runs) & \code{fortran\_\allowbreak{}diff\_\allowbreak{}D057\_\allowbreak{}rho\_\allowbreak{}kernels}, \code{l1\_\allowbreak{}peerho} \\
66 & D-057 F21 rho NaN pion & rho\_\allowbreak{}decay.\allowbreak{}f:\allowbreak{}111-112 & acos(pz/\allowbreak{}p/\allowbreak{}cos θ\_oop) with an argument above 1 by rounding gives th\_inp = NaN, yptar = NaN; mc\_shms and pass\_cuts let it through and the event contributes with every recon pion column NaN (1 per 1e5 events: rho\_\allowbreak{}h\_\allowbreak{}k1\_\allowbreak{}rad\_\allowbreak{}A record 57258, rho\_\allowbreak{}he3\_\allowbreak{}k1\_\allowbreak{}rad\_\allowbreak{}A 54227) & the argument is clamped to [−1, 1] (\code{rho::\allowbreak{}Rho\allowbreak{}Decay\allowbreak{}Config::\allowbreak{}acos\_\allowbreak{}clamp}\simcidx{Classes and interfaces!RhoDecayConfig@\code{RhoDecayConfig}}) & Fortran (original behaviour) & about 1 event per 1e5 (its weight \textasciitilde{}6e-8 of the sum); NaN tree entries disappear & \code{fortran\_\allowbreak{}diff\_\allowbreak{}D057\_\allowbreak{}rho\_\allowbreak{}kernels} \\
67 & D-057 F15 rho .hist title & simc.\allowbreak{}f:\allowbreak{}720-725 & the .hist header says 'I am not set up for anything else yet!' for D and ³He rho runs, which run & \code{\ ****-{}-{}-{}-{}-{}-{}-{}-{}\ \ D(\allowbreak{}e,\allowbreak{}e'rho)\ \ -{}-{}-{}-{}-{}-{}-{}-{}****} / \code{A(\allowbreak{}e,\allowbreak{}e'rho)} (\code{RunReport}\simcidx{Classes and interfaces!RunReport@\code{RunReport}}) & Fortran (original behaviour) & .hist header of D/³He rho runs & \code{fortran\_\allowbreak{}diff\_\allowbreak{}D057\_\allowbreak{}F15\_\allowbreak{}F18\_\allowbreak{}rho} \\
68 & D-057 F18 rho P-arm report & simc.\allowbreak{}f:\allowbreak{}259-261,\allowbreak{} 283-285, 316-317, 324-325; event.\allowbreak{}f:\allowbreak{}22-24 & the .geni/.gen P-arm histograms, the .hist 'Hadron arm' resolution sums and the 'Limiting VERTEX P arm' ranges use vertex\%p, which for doing\_rho is the RHO with its lab theta/phi in the slope fields (rho\_\allowbreak{}h\_\allowbreak{}k1\_\allowbreak{}rad: Ave.Error yptar −1580.9 mr, limits 589…2335 mr) & the decay pion orig\%p (\code{Generated\allowbreak{}Hadron\allowbreak{}Arm}, corrected mode) & Fortran (original behaviour) & .geni, .gen (and gen.root) P-arm distributions,\allowbreak{} .hist resolution lines and limits; no weight & \code{fortran\_\allowbreak{}diff\_\allowbreak{}D057\_\allowbreak{}F15\_\allowbreak{}F18\_\allowbreak{}rho} \\
69 & U-A1 nuclear elastic deleted & upstream 59ccdf6 event.\allowbreak{}f:\allowbreak{}507/\allowbreak{}517/\allowbreak{}546/\allowbreak{}937/\allowbreak{}1344/\allowbreak{}1399/\allowbreak{}1457,\allowbreak{} NtupleInit.\allowbreak{}f:\allowbreak{}293,\allowbreak{} results\_\allowbreak{}write.\allowbreak{}f:\allowbreak{}246 (removed by e4dd175, added by 4df7351) & e4dd175 rebuilt event.\allowbreak{}f/\allowbreak{}NtupleInit.\allowbreak{}f/\allowbreak{}results\_\allowbreak{}write.\allowbreak{}f from a pre-4df7351 copy: \code{doing\_\allowbreak{}nuc\_\allowbreak{}elast\ =\ 1} (shipped infiles/\allowbreak{}ee\_\allowbreak{}calcium48.inp) leaves vertex\%\allowbreak{}e\%\allowbreak{}E = 0, every try fails in gen\_rad (\code{Negative\ electron\ energy}), 0 successes, the run never ends & the complete A(e,e')A path as at commit 20d46fc (reaction \code{nuclear\_\allowbreak{}elastic}, model \code{nuc\_\allowbreak{}elastic}, the (e,e'p) ntuple layout), in both modes: \code{-{}-{}fortran-{}compat}\simcidx{Command-line options!fortran-compat@\code{-{}-{}fortran-{}compat}} cannot reproduce an endless run and follows 20d46fc (task U.1) & Fortran (original behaviour) & doing\_\allowbreak{}nuc\_\allowbreak{}elast runs only. ee\_\allowbreak{}calcium48 (5000 events): .hist/\allowbreak{}.geni/\allowbreak{}.gen/\allowbreak{}h10 identical to the 20d46fc Fortran (reference/\allowbreak{}runs/\allowbreak{}up\_\allowbreak{}ee\_\allowbreak{}calcium48\_\allowbreak{}20d46fc) with RANLUX + \code{-{}-{}fortran-{}compat}\simcidx{Command-line options!fortran-compat@\code{-{}-{}fortran-{}compat}} & \code{l1\_\allowbreak{}nuc\_\allowbreak{}elastic}, \code{l1\_\allowbreak{}simcrun\_\allowbreak{}ranlux\_\allowbreak{}up\_\allowbreak{}ee\_\allowbreak{}calcium48\_\allowbreak{}cap1000\_\allowbreak{}20d46fc}, \code{l1\_\allowbreak{}simcrun\_\allowbreak{}ranlux\_\allowbreak{}up\_\allowbreak{}ee\_\allowbreak{}calcium48\_\allowbreak{}20d46fc} \\
70 & U-A2 kaon Nntu STOP & upstream 59ccdf6 results\_\allowbreak{}write.\allowbreak{}f:\allowbreak{}222-228 vs NtupleInit.\allowbreak{}f (kaon branch) & every kaon run with Nntu = 1 STOPs at the first event (\code{number\ of\ output\ quantities\ not\ the\ same\ as\ output\ array\ size}: 66 values written, 65 tags) & corrected in both modes (U.2, D-082): the upstream layout (the SIMC++ default) declares every written column, \code{sigcm\ saghai\ factor} (66 columns; src/\allowbreak{}output/\allowbreak{}NtupleLayout.\allowbreak{}cpp UpstreamMesonColumns); the legacy 294f662 layout has 55 & Fortran (original behaviour) & kaon runs with Nntu = 1 (upstream 59ccdf6 only) & fortran\_\allowbreak{}diff\_\allowbreak{}U\_\allowbreak{}A2\_\allowbreak{}A4\_\allowbreak{}upstream\_\allowbreak{}layout\_\allowbreak{}kaon\_\allowbreak{}pi0\_\allowbreak{}field \\
71 & U-A4 π0 + field ntuple overflow & upstream 59ccdf6 results\_\allowbreak{}write.\allowbreak{}f:\allowbreak{}10 \code{ntu(\allowbreak{}80)}, NtupleInit.\allowbreak{}f:\allowbreak{}11 \code{NtupleTag(\allowbreak{}80)}, util/\allowbreak{}root\_\allowbreak{}tree make\_\allowbreak{}root\_\allowbreak{}tree arrays & π0 with using\_\allowbreak{}tgt\_\allowbreak{}field writes 84 values into arrays of 80 (out of bounds; latent, never run) & SIMC++ sizes the ntuple from the layout (no fixed array), in both modes: the upstream layout (default) writes all 84 columns (photons 65-76, field angles 77-84); the legacy 294f662 layout has 73 columns (D-067 T14) & Fortran (original behaviour) & none in SIMC++ & fortran\_\allowbreak{}diff\_\allowbreak{}U\_\allowbreak{}A2\_\allowbreak{}A4\_\allowbreak{}upstream\_\allowbreak{}layout\_\allowbreak{}kaon\_\allowbreak{}pi0\_\allowbreak{}field \\
\end{simctable}
\end{landscape}
\subsection*{\texorpdfstring{Category B: ambiguous, kept (the expert list)}{Category B: ambiguous, kept (the expert list)}}\label{translation_report-06u5fcatalogue:category-b-ambiguous-kept-the-expert-list}
\begin{landscape}
\begin{simctable}{scriptsize}{8}{l>{\raggedright\arraybackslash}p{\dimexpr 0.104\linewidth-2\tabcolsep\relax}>{\raggedright\arraybackslash}p{\dimexpr 0.099\linewidth-2\tabcolsep\relax}>{\raggedright\arraybackslash}p{\dimexpr 0.213\linewidth-2\tabcolsep\relax}>{\raggedright\arraybackslash}p{\dimexpr 0.099\linewidth-2\tabcolsep\relax}>{\raggedright\arraybackslash}p{\dimexpr 0.102\linewidth-2\tabcolsep\relax}>{\raggedright\arraybackslash}p{\dimexpr 0.123\linewidth-2\tabcolsep\relax}>{\raggedright\arraybackslash}p{\dimexpr 0.136\linewidth-2\tabcolsep\relax}}
\caption{Category B: ambiguous, kept (the expert list)}\\%
\toprule\rowcolor{simchead}\simcth{\#} & \simcth{id} & \simcth{Fortran location} & \simcth{Fortran behaviour} & \simcth{default (corrected) mode} & \simcth{\code{-{}-{}fortran-{}compat}} & \simcth{effect} & \simcth{question / marker} \\\midrule\endfirsthead
\multicolumn{8}{@{}l}{\simccontinued}\\\toprule\rowcolor{simchead}\simcth{\#} & \simcth{id} & \simcth{Fortran location} & \simcth{Fortran behaviour} & \simcth{default (corrected) mode} & \simcth{\code{-{}-{}fortran-{}compat}} & \simcth{effect} & \simcth{question / marker} \\\midrule\endhead
\midrule\multicolumn{8}{r@{}}{\simccontinues}\\\endfoot
\bottomrule\endlastfoot
72 & D-033 F4 D jacobian & event.\allowbreak{}f:\allowbreak{}1020-1024 & D(e,e'p): the hadron 1/\allowbreak{}cos³(θ−θ0) jacobian is applied only for doing\_\allowbreak{}heavy/\allowbreak{}pion/\allowbreak{}kaon/\allowbreak{}delta/\allowbreak{}semi,\allowbreak{} not for D, although the proton angles are generated in xptar/yptar and genvol contains dΩp (simc.\allowbreak{}f:\allowbreak{}388) & Fortran kept & Fortran kept & D(e,e'p) weights: the factor would be 1/cos³ ≈ 1 − 1.5(xptar²+yptar²),\allowbreak{} i.e. 0 at the centre to −1.5 \% at the corners of the generated hadron solid angle (±60 × ±80 mr) & Question: is the missing jacobian for D intended? \\
73 & D-033 F1 SF fallback & init.\allowbreak{}f:\allowbreak{}839-851,\allowbreak{} dbase.\allowbreak{}f:\allowbreak{}598-616 & every A other than 2/12/56/197 silently uses c12.theory (benharsf\_\allowbreak{}12.dat with use\_\allowbreak{}benhar\_\allowbreak{}sf): He-3, He-4, Ca-40 get the C-12 spectral function. SIMC++ keeps the default but prints a WARNING and lets the input name the file (\code{pipeline.\allowbreak{}reaction.\allowbreak{}spectral\_\allowbreak{}function.\allowbreak{}file}) & Fortran kept & Fortran kept & eep\_\allowbreak{}he4\_\allowbreak{}rad\_\allowbreak{}A is "He-4 kinematics with the C-12 SF" & Question: which spectral function for He-3/He-4 (theory files exist only for 2/\allowbreak{}12/\allowbreak{}56/\allowbreak{}197)? \\
74 & D-015 F7 HMS hodo 2 & hms/\allowbreak{}mc\_\allowbreak{}hms\_\allowbreak{}hut.\allowbreak{}f:\allowbreak{}537-551 & hodoscope-2 planes tested with the hodoscope-1 limits/\allowbreak{}offsets & Fortran kept & Fortran kept & HMS acceptance at the hut (scin trigger) & Question: copy-paste or the same active area on purpose? \\
75 & D-015 F7 HMS Cherenkov & hms/\allowbreak{}mc\_\allowbreak{}hms\_\allowbreak{}hut.\allowbreak{}f:\allowbreak{}523-529 & multiple scattering in the Cherenkov exit foil applied twice & Fortran kept & Fortran kept & HMS fp smearing (\code{mc\_\allowbreak{}smear}\simcidx{Parameters!mc_smear@\code{mc\_smear}}) & Question: double counting or a second foil? \\
76 & D-015 F5 SHMS drifts & shms/\allowbreak{}mc\_\allowbreak{}shms.\allowbreak{}f (zd\_d1men 39.47, zd\_d1out 60.68, zd\_fp 307.95) & hard-coded drift lengths differ from the COSY classes (39.4715, 60.68058, 307.9514) & Fortran kept & Fortran kept & SHMS fp ≤ 2.2e-4 cm, path −0.0035 cm & Question: use the COSY values? \\
77 & D-017 can radius & target.\allowbreak{}f:\allowbreak{}368-376,\allowbreak{} 456-464 vs 108, 228 & beer-can corner of \code{extreme\_\allowbreak{}trip\_\allowbreak{}thru\_\allowbreak{}target} uses 1.25", \code{trip\_\allowbreak{}thru\_\allowbreak{}target} 1.325" & Fortran kept & Fortran kept & extreme Eloss/teff limits (beer can) & Question: which radius? \\
78 & D-012 F4 th\_can & target.\allowbreak{}f:\allowbreak{}236-244 & for |z\_can/R| > 1 + 1e-6 the pudding-can \code{th\_\allowbreak{}can} keeps its previous value (warning only) & Fortran kept & Fortran kept & can 2 edge events & Question: reject, or clamp? \\
79 & D-019 one\_tail & radc.\allowbreak{}f:\allowbreak{}198-216,\allowbreak{} 234, 358 & with \code{one\_\allowbreak{}tail\ ≠\ 0} a disabled tail can be drawn: no photon, gen\_weight = 1/\allowbreak{}hardcorfac & Fortran kept & Fortran kept & runs with one\_tail ≠ 0 & Question: restrict \code{frac} to the enabled tails? \\
80 & D-019 doing\_\allowbreak{}proton & init.f:693 & \code{doing\_\allowbreak{}tail(\allowbreak{}3)} passed as \code{doing\_\allowbreak{}proton} to lambda\_\allowbreak{}dave: switching off tail 3 removes the ep interference from lambda(1:2) & Fortran kept & Fortran kept & one\_tail runs & Question: intended? \\
81 & D-019 intcor\_\allowbreak{}mode 0 & init.\allowbreak{}f:\allowbreak{}701-715 & \code{schwinger} does not set \code{dsoft\_\allowbreak{}prime}; the SAVEd value (0) is used, g(4) loses its internal part (0.019443 vs 0.098637) & Fortran kept & Fortran kept & intcor\_\allowbreak{}mode = 0 runs & Question: derivative term intended? \\
82 & D-019 rad\_weight & radc.\allowbreak{}f:\allowbreak{}330-333,\allowbreak{} 403-406, 454-457 & \code{rad\_\allowbreak{}weight*basicrad\_\allowbreak{}weight} overwritten by \code{peaked\_\allowbreak{}rad\_\allowbreak{}weight} (rad\_flag ≤ 1) & Fortran kept & Fortran kept & radiative weights & Question: intended (peaked weight contains basicrad)? \\
83 & D-019 tail-1 COMMON & radc.\allowbreak{}f:\allowbreak{}322-328,\allowbreak{} 512 & the \code{complete\_\allowbreak{}ev} inside tail 1 recomputes hardcorfac/\allowbreak{}frac/\allowbreak{}g/\allowbreak{}c before the final division & Fortran kept & Fortran kept & radiative weights & Question: post-radiation hardcorfac intended? \\
84 & D-019 Q4 phi\_ext & radc.\allowbreak{}f:\allowbreak{}566-642 & \code{rad\_\allowbreak{}flag\ =\ 0} peaked weight independent of Egamma; phi\_ext, dsoft\_\allowbreak{}ext*,\allowbreak{} Friedrich terms unused & Fortran kept & Fortran kept & rad\_flag = 0 & Question: include phi\_ext? \\
85 & D-019 Q4 dsoft\_int & radc.\allowbreak{}f:\allowbreak{}590-620 & rad\_flag = 1 with extrad\_\allowbreak{}flag = 3 uses the SAVEd dsoft\_\allowbreak{}intmin/\allowbreak{}max (0) & Fortran kept & Fortran kept & that flag combination & Question: compute them? \\
86 & D-019 Egamma\_\allowbreak{}min(1) & radc.\allowbreak{}f:\allowbreak{}283-286 & for ntail = 0 \code{Egamma\_\allowbreak{}min(\allowbreak{}1)} keeps its stale COMMON value & Fortran kept & Fortran kept & D(e,e'p) rad\_flag 3 (P8) & Question: set it? \\
87 & D-005 lambda input & init.\allowbreak{}f:\allowbreak{}692-695 & the input \code{lambda(\allowbreak{}1:\allowbreak{}3)} is overwritten every event & Fortran kept & Fortran kept & none (input ignored) & Question: drop the input or honour it? \\
88 & D-024 stale orig & event.\allowbreak{}f:\allowbreak{}400-415 & failed \code{generate\_\allowbreak{}rad} leaves \code{orig} from the previous try & Fortran kept & Fortran kept & per-try outputs only & Question: — (try rejected) \\
89 & D-024 grnd dEbeam=0 & event.f:204 & a grnd is drawn for Ein even when dEbeam = 0 & Fortran kept & Fortran kept & random stream only & Question: — \\
90 & D-024 Eloss gaussians & event.\allowbreak{}f:\allowbreak{}193-205,\allowbreak{} trip\_\allowbreak{}thru\_\allowbreak{}target & energy-loss gaussians drawn even with using\_\allowbreak{}Eloss = 0 & Fortran kept & Fortran kept & random stream only & Question: — \\
91 & D-015 F2 HMS resmult & hms/\allowbreak{}mc\_\allowbreak{}hms\_\allowbreak{}hut.\allowbreak{}f:\allowbreak{}292 & one grnd per hut entry even with smearing off & Fortran kept & Fortran kept & random stream only & Question: — \\
92 & D-015 tmpfact & simc.f:1334 & SAVEd \code{tmpfact}: when the electron arm stops before the hut, \code{ntup\%\allowbreak{}resfac} gets the previous event's value & Fortran kept & Fortran kept & ntup resfac & Question: intended? \\
93 & D-015 F6 SHMS counters & shms/\allowbreak{}mc\_\allowbreak{}shms\_\allowbreak{}hut.\allowbreak{}f & S1 x and y both count in s1; 2x counts in s3, 2y in s2; \code{cal} counts two cuts (true position, fitted track 5 cm inside) & Fortran kept & Fortran kept & .geni diagnostics & Question: separate counters? \\
94 & D-011 recon xptar & simc.f:274, 277 vs 256, 262 & the .gen \code{recon} xptar histograms are filled NOT negated, \code{gen/\allowbreak{}geni} with −vertex xptar & Fortran kept & Fortran kept & .gen / .gen.root histograms & Question: which convention? \\
95 & D-015 F4 fry & simc.\allowbreak{}f:\allowbreak{}1467,\allowbreak{} 1487, 1733, 1753 vs 1472-1482 & mc\_\allowbreak{}hms/\allowbreak{}mc\_\allowbreak{}shms get \code{xtar\_\allowbreak{}init} as fry, SOS/HRS get the raster fry & Fortran kept & Fortran kept & HMS/SHMS recon (fry terms) & Question: intended? \\
96 & D-024 H elastic W & event.f:723 block (pion/kaon only) & complete\_\allowbreak{}ev never sets vertex W, epsilon, theta\_pq, phi\_pq for H elastic & Fortran kept & Fortran kept & vertex-only quantities & Question: intended? \\
97 & D-024 ngen<0 & simc.\allowbreak{}f:\allowbreak{}169-199,\allowbreak{} 341-350 & with ngen < 0 the .hist wtcontr and AVERAGE.sigcc are divided by tries & Fortran kept & Fortran kept & .hist with ngen < 0 & Question: intended? \\
98 & D-018 herho & dbase.f:213 & \code{doing\_\allowbreak{}herho\ =\ A.\allowbreak{}eq.\allowbreak{}3} (the other he* flags use .ge.3) & Fortran kept & Fortran kept & rho on A > 3 (P13) & Question: intended? \\
99 & D-005 sigc\_nbin & semi\_\allowbreak{}physics.\allowbreak{}f:\allowbreak{}173 & default \code{sigc\_\allowbreak{}nbin\ =\ 0} → division by zero & Fortran kept & Fortran kept & SIDIS (P12) & Question: required input? \\
100 & D-038 F4 coherent model & physics\_\allowbreak{}pion.\allowbreak{}f:\allowbreak{}807-854 (exclfit), event.\allowbreak{}f:\allowbreak{}1496-1500 & coherent pion production (which\_\allowbreak{}pion 10/11) uses the nucleon fit with Mtar\_\allowbreak{}struck = M\_A (W = √s photon-nucleus 4–13 GeV); exclfit keeps \code{mtar\_\allowbreak{}gev\ =\ 0.\allowbreak{}938} (l.812); tgtweight = Z or N ("OK for coherent???"); sig\_\allowbreak{}param\_\allowbreak{}2021 is negative in 210 of 4564 table rows (no clamping) & Fortran kept & Fortran kept & coherent pion weights (P9) & Question: is a nucleon model meant for coherent production, and should negative cross sections be clamped? \\
101 & D-038 F3 coherent Mrec & dbase.\allowbreak{}f:\allowbreak{}283-288,\allowbreak{} 380-388 & for coherent production the input \code{targ\%\allowbreak{}mrec\_\allowbreak{}amu}\simcidx{Parameters!targ.mrec_amu@\code{targ\%mrec\_amu}} is first replaced by the A−1 guess M − Mp, then by M − M\_\allowbreak{}N(struck) + M\_\allowbreak{}N(recoil); the input value is never used, even when it is the correct final nucleus (two replacement messages) & Fortran kept & Fortran kept & coherent pion runs: Mrec\_\allowbreak{}struck (He-3 coh: 2809.68 MeV) & Question: should the input final-nucleus mass be honoured? \\
102 & D-038 F9 Δ re-generation & event.\allowbreak{}f:\allowbreak{}620-631,\allowbreak{} radc.f:324 & complete\_\allowbreak{}ev overwrites the COMMON targ\%\allowbreak{}Mrec\_\allowbreak{}struck with a new Breit-Wigner Δ mass on every call; radiation tail 1 calls complete\_\allowbreak{}ev again, so the event gets a second Δ mass and (π0) a second photon decay; the \code{.\allowbreak{}hist} prints the last Δ mass of the run (simc.\allowbreak{}f:\allowbreak{}844; SIMC++ too, \code{IReaction::\allowbreak{}final\allowbreak{}Target\allowbreak{}State}\simcidx{Classes and interfaces!IReaction@\code{IReaction}}) & Fortran kept & Fortran kept & π Δ final states (which\_\allowbreak{}pion 2/3) with radiation & Question: intended re-generation? \\
103 & D-038 F12 sig\_\allowbreak{}multipole & physics\_\allowbreak{}pion.\allowbreak{}f:\allowbreak{}197-398 (dead code: the call in peepi is commented out, l.157-163) & COMPLEX*8 (single-precision) amplitudes although the build is -fdefault-real-8; at exact (Q², ν) nodes or outside the table the amplitudes of the PREVIOUS call are used; the argument \code{nu} is overwritten (≥ 500 → 499.9); 197.327/mh fixed at the first call. SIMC++ reproduces all of it (\code{pion::\allowbreak{}Sig\allowbreak{}Multipole}\simcidx{Classes and interfaces!SigMultipole@\code{SigMultipole}}, \code{std::\allowbreak{}complex<float>} storage, Fortran kind rules per operation, SAVEd members) & Fortran kept & Fortran kept & none (not called by SIMC); L1 sig\_\allowbreak{}multipole 612/612 rows bitwise & Question: single precision and the previous-call amplitudes intended? \\
104 & D-037 F3 coll scattering & hms/\allowbreak{}mc\_\allowbreak{}hms\_\allowbreak{}coll.\allowbreak{}f:\allowbreak{}115-118,\allowbreak{} shms/\allowbreak{}mc\_\allowbreak{}shms\_\allowbreak{}coll.\allowbreak{}f:\allowbreak{}117-120 & musc and the Landau-sampled enerloss\_\allowbreak{}new (2 + 1 gauss1) of every collimator step are applied whatever ms\_flag / mc\_smear say & Fortran kept & Fortran kept & pion/muon runs with using\_\allowbreak{}*coll and mc\_smear = 0 & Question: is the collimator scattering intended physics (the purpose of the stepping) or should it follow ms\_flag? \\
105 & D-037 F5 pion\_\allowbreak{}coll\_\allowbreak{}absorb & hms/\allowbreak{}pion\_\allowbreak{}coll\_\allowbreak{}absorb.\allowbreak{}f:\allowbreak{}16,\allowbreak{} 30-56 & own mpi = 139.56995 (constants.\allowbreak{}inc 139.57018); transmission 1 for Tpi ≤ 85 MeV, 0.938 per step just above; unused mate\_t; the SHMS uses the HMS material (17 g/cm³, A = 171.57) & Fortran kept & Fortran kept & absorption probability of pions in both collimators & Question: intended parameters? extrapolate below 85 MeV? \\
106 & D-037 F8 transp decay & shared/\allowbreak{}transp.\allowbreak{}f:\allowbreak{}142-185,\allowbreak{} 232-276 & a first-half decay is applied at the element entrance (then the map acts on the daughter), a second-half decay at the exit with the parent's β, γ; decdist gets z\_decay, pathlen zd+δz & Fortran kept & Fortran kept & decays inside magnets (doing\_\allowbreak{}decay) & Question: an intended approximation (the COSY map cannot be split)? \\
107 & D-043 F2 data & saghai\_\allowbreak{}sigma0.dat (content) & the file starts with 18 stray records (s = 2.898 GeV², Q² = 0.52/0.75) and has 1 record instead of 19 for s = 2.851, Q² = 0; line 72 has a truncated field (\code{-{}00E-{}02}, read as −0) & Fortran kept & Fortran kept & \code{saghai} column near the Σ threshold at Q² < 0.25 GeV² only & Question: which amplitudes belong to s = 2.851, Q² = 0? \\
108 & D-043 F7 decay model & shared/\allowbreak{}transp.\allowbreak{}f:\allowbreak{}149-160,\allowbreak{} 239-251; project.\allowbreak{}f:\allowbreak{}76-85 & K decay: fixed 70 \% μν / 30 \% ππ, p* = sqrt(Mk²/4 − Mπ²) = 203.6 MeV/c (equal-mass formula with the charged pion; K\textsuperscript{+}→π\textsuperscript{+}π\textsuperscript{0} is 205.1), the daughters never decay & Fortran kept & Fortran kept & decay-ON kaon runs (the Mk of transp.\allowbreak{}f:\allowbreak{}158 is corrected, D-006) & Question: intended approximation? \\
109 & D-043 F8 Saghai angles & physics\_\allowbreak{}kaon.\allowbreak{}f:\allowbreak{}19,\allowbreak{} 79-85 & peeK's \code{phi} is never assigned (SAVEd, 0) and eekeek's \code{theta} argument is main\%\allowbreak{}theta\_\allowbreak{}pq (the lab p–q angle) & Fortran kept & Fortran kept & \code{saghai} column & Question: intended (comment "ALWAYS CALCULATING FOR PHI=0")? \\
110 & D-043 F9 coherent kaon & physics\_\allowbreak{}kaon.\allowbreak{}f:\allowbreak{}82-86,\allowbreak{} 215; event.\allowbreak{}f:\allowbreak{}1505 & coherent production (which\_\allowbreak{}kaon 10/11/12) uses the nucleon formulas with W of the photon–nucleus system (4–13.5 GeV); with Mrec\_\allowbreak{}struck > 1150 MeV the coherent Λ uses the Σ branches; the cross section is the same for 10/11/12 (only tgtweight Z/N differs, "OK for coherent???") & Fortran kept & Fortran kept & coherent kaon weights & Question: is a nucleon model meant for coherent production? \\
111 & D-043 F10 Mp & physics\_\allowbreak{}kaon.\allowbreak{}f:\allowbreak{}194-197 & sig\_\allowbreak{}factorized uses the proton mass in ν/q\_cm also for Σ\textsuperscript{−} (neutron) and coherent production & Fortran kept & Fortran kept & Σ\textsuperscript{−} and coherent weights & Question: intended? \\
112 & \textbf{D-048 F4 Delta weight} & physics\_\allowbreak{}delta.\allowbreak{}f:\allowbreak{}130-132,\allowbreak{} event.\allowbreak{}f:\allowbreak{}1512 & \textbf{peedelta = 1.0\emph{jacobian}(gtpr*fac): the cross-section line is commented out ("Need new xsec model."), so the Delta weight is phase space × photon flux × the transform\_\allowbreak{}to\_\allowbreak{}cm jacobian, with NO cross section.} The sig\_blok value it computes (ntup\%\allowbreak{}sigcm) is the Brauel p(e,e'π+)n fit evaluated with the proton's CM angles and the proton t (π− branch for which\_pion 1/3/11): not a pπ0 or Δ model either & Fortran kept & Fortran kept & every Delta weight, wtcontr, \code{Weight}, \code{siglab} (14.8–351 in delta\_\allowbreak{}h\_\allowbreak{}d1\_\allowbreak{}rad\_\allowbreak{}A) & Question: should SIMC++ offer a real Delta (e,e'p)π cross section, and which one? (no choice can be derived from the Fortran) \\
113 & D-048 F5 two roots & event.\allowbreak{}f:\allowbreak{}682-684 & for doing\_\allowbreak{}delta one grnd picks the second root of the proton-energy quadratic when grnd() > 0.5; no factor 1/0.5 enters gen\_weight or the jacobian, so when both roots are physical each is generated at half rate with full weight (and the first-root check \code{p\%\allowbreak{}E\ <\ 0} comes before the choice) & Fortran kept & Fortran kept & Delta weights where both roots reach the acceptance (D2: 3 of 1000 capture events took the second root; D1: none) & Question: is a factor 2 missing? \\
114 & D-048 F5b root re-choice & radc.f:322 with event.\allowbreak{}f:\allowbreak{}682-684 & with tail 1 radiating (corrected F2), generate\_\allowbreak{}rad calls complete\_\allowbreak{}ev again, which draws a NEW root choice (the radiated event may take the other root than the one the tail limits were computed with), as D-038 F9 for the Δ mass & Fortran kept & Fortran kept & corrected-mode rad Delta runs (tail-1 events) & Question: intended re-generation? \\
115 & D-048 F7 recoil mass & dbase.\allowbreak{}f:\allowbreak{}297-300 & targ\%\allowbreak{}Mrec\_\allowbreak{}struck = Mpi = 139.57018 (charged pion) although for γ*p → p X the missing pion is a π0 (134.9766) & Fortran kept & Fortran kept & generated Delta kinematics; central recon MM 138.92 MeV & Question: π0 mass intended? \\
116 & D-048 F8 which\_pion & event.\allowbreak{}f:\allowbreak{}470-491,\allowbreak{} 612-623; physics\_\allowbreak{}delta.\allowbreak{}f:\allowbreak{}99 & complete\_\allowbreak{}ev reads which\_pion also for doing\_\allowbreak{}delta: with 2/3 it reads delta\_\allowbreak{}relativistic\_\allowbreak{}bw.\allowbreak{}inp and replaces the recoil pion by a Breit-Wigner Δ mass on every call (below threshold in the Δ region: every complete\_\allowbreak{}ev fails, an L2 run would never finish); 1/3/11 select the π− branch of sig\_blok (only \code{sigcm} changes) & Fortran kept & Fortran kept & Delta runs with which\_pion ≠ 0 & Question: should which\_pion be ignored for doing\_\allowbreak{}delta? \\
117 & D-051 F7 threshold & semi\_\allowbreak{}physics.\allowbreak{}f:\allowbreak{}278-284 & the "2π / πK threshold" cut uses mmpi2 = M² + 2Mν(1−z)(1−pt²[GeV²]) — a mass² multiplied by (1 − pt²) — and returns 0 below (M + m\_h)² & Fortran kept & Fortran kept & weight 0 or not for events at large z / large pt² & Question: which missing-mass formula was intended? \\
118 & D-051 F8 semi thresholds & event.\allowbreak{}f:\allowbreak{}979-996 & complete\_\allowbreak{}ev rejects a SIDIS try when (Mtar\_\allowbreak{}struck + ν − E\_h)² − P\_miss² < (Mp + Mpi0)² — for every hadron (also kaons, where the lowest final state is ΛK), with Mtar\_\allowbreak{}struck and no Fermi motion; the following z > 1 test is unreachable (z > 1 already fails the threshold); 71 \% of all tries are z > 1 events & Fortran kept & Fortran kept & generation efficiency and the accepted phase space near z → 1 (kept in both modes, \code{Sidis\allowbreak{}Production::\allowbreak{}complete\allowbreak{}Ev}\simcidx{Classes and interfaces!SidisProduction@\code{SidisProduction}}) & Question: which threshold is intended for kaons / with Fermi motion? \\
119 & D-051 F16 hadron range & init.\allowbreak{}f:\allowbreak{}504-506 & with radiation gen\%p\%E reaches sumEgen\%\allowbreak{}max − edge\%\allowbreak{}e\%\allowbreak{}E\%\allowbreak{}min (hadron δ up to +116 \%); 2.9 \% of the rad events have orig δ > 80 \% and 1612/100000 of them reconstruct inside −10…22 \% (SHMS model) & Fortran kept & Fortran kept & 0.68 \% efficiency (rad) vs 4.8 \% (norad); 1.2 \% of ΣWeight from δ > 80 \% (kept, set-up) & Question: intended generation range? \\
120 & D-051 F10 model choices (R\_h for D) & semi\_\allowbreak{}physics.\allowbreak{}f:\allowbreak{}659-677 & Rhad\_\allowbreak{}global (the nuclear hadron-attenuation fit) returns 1 only for A = 1: deuterium gets N\textsubscript{0} z\textasciicircum{}α (1−z)\textasciicircum{}β,\allowbreak{} 0.955 at z = 0.5 & Fortran kept & Fortran kept & D(e,e′h)X weights (≈ −4.5 \%) & Question: intended for D? (D-051 F10 also lists the slacemcfit and Mtar\_\allowbreak{}struck = Mp choices) \\
121 & D-043 F12 warning & physics\_\allowbreak{}kaon.\allowbreak{}f:\allowbreak{}201-206 & the "Kaon C.M. momentum passed to sig\_\allowbreak{}factorized does not agree" warning is printed for the central event of Σ\textsuperscript{0},\allowbreak{} Σ\textsuperscript{−} and A > 1 quasifree runs (vertex0 places the kaon at spec\%p\%P whatever the hyperon) & Fortran kept & Fortran kept & log text & Question: message only \\
122 & D-057 F4 rho nuclei & dbase.\allowbreak{}f:\allowbreak{}563,\allowbreak{} 594; event.\allowbreak{}f:\allowbreak{}336-372; sf\_\allowbreak{}lookup.\allowbreak{}f:\allowbreak{}195-245 & no momentum distribution (nump = 0) and no ³He spectral function are loaded for doing\_\allowbreak{}deutrho/\allowbreak{}herho: the Fermi block draws 4 uniforms but pfer = 0 always; ³He generate\_\allowbreak{}em reads its empty table out of bounds and returns Em = 0 (2 uniforms); D: Em = 2.22494, efer = Md − Mn & Fortran kept & Fortran kept & D/³He rho runs are a nucleon at rest (rho\_\allowbreak{}d\_\allowbreak{}k1\_\allowbreak{}rad\_\allowbreak{}A: wtcontr = 2.33 × H, i.e. Z + N with the same kinematics) & Question: should D/³He rho include Fermi motion / a spectral function? (a model limitation, kept in both modes; SIMC++ computes the ³He Em = 0 without the out-of-bounds read) \\
123 & D-057 F5 rho s & generate\_\allowbreak{}rho.\allowbreak{}f:\allowbreak{}72 & s = −Q² + efer² + 2ν·efer lacks the −2q·pfer − pfer² terms while the boost includes them & Fortran kept & Fortran kept & none today & Question: intended? (latent: pfer = 0 always, F4) \\
124 & D-057 F7 rho t & rho\_\allowbreak{}physics.\allowbreak{}f:\allowbreak{}115,\allowbreak{} 233 & t and tmin of peerho use the nominal Mrho², not the event's Breit-Wigner mass (complete\_\allowbreak{}ev's main\%t uses it) & Fortran kept & Fortran kept & the exp(−b· & Question: model defined at the pole mass on purpose? \\
125 & D-057 F8 rho decay frame & rho\_\allowbreak{}decay.\allowbreak{}f:\allowbreak{}57-77 & the decay angles (cos θ* from sin²θ* + 2εR cos²θ\emph{, φ} flat) are thrown w.r.t. the LAB z/x axes in the rho rest frame, not the helicity frame; ntup\%\allowbreak{}rhotheta is that angle. (The header comment says εR, the code 2εR = SCHC; the envelope is valid) & Fortran kept & Fortran kept & the detected pion's angular distribution (the 2εR cos² part) & Question: intended approximation? \\
126 & D-057 F10 rho xptar & rho\_\allowbreak{}decay.\allowbreak{}f:\allowbreak{}110,\allowbreak{} 142 & xptar = tan(asin(px/\allowbreak{}p)) instead of px/pz′ (up to 4.3e-4 inside ±0.1; orig\%\allowbreak{}p\%\allowbreak{}theta/\allowbreak{}phi differ from orig\%up by ≤ 0.42 mr) & Fortran kept & Fortran kept & decay-pion xptar & Question: approximation accepted? \\
127 & D-057 F11 rho tail 3 & radc.\allowbreak{}f:\allowbreak{}386-387; init.\allowbreak{}f:\allowbreak{}231-232 & tail-3 radiation of the NEUTRAL rho, with limits from the hadron edges computed with Mrho: 70 \% of tail-3 tries fail at the limits (91682 / 131589 H tries) & Fortran kept & Fortran kept & rad efficiency (rad/norad wtcontr 0.854); accepted events have Egamma(3) up to 1.5 GeV & Question: should the rho radiate, and with which limits? \\
128 & D-057 F12 rho tail 1 & radc.f:324 & a tail-1 photon re-calls complete\_\allowbreak{}ev,\allowbreak{} which generates a NEW rho (new BW mass and CM angles; as D-038 F9) & Fortran kept & Fortran kept & rad events of tail 1 (393 of 1000 capture events) & Question: intended? \\
129 & D-012 F4 calo windows & target.\allowbreak{}f:\allowbreak{}73-102,\allowbreak{} 189-222 & no window branch for the calorimeter arms 7/8: s\_air, s\_kevlar, s\_mylar and forward\_\allowbreak{}path keep the previous narm 2/3 call's values (static locals; the first ones from init.\allowbreak{}f:\allowbreak{}53-56 target\_\allowbreak{}init),\allowbreak{} s\_Al = 0 & Fortran kept & Fortran kept & energy loss and radiation length of the calorimeter-arm particle (π0 photons via the π0, the calorimeter electron) & Question: which windows / target path should the calorimeter arm get (NPS geometry)? kept in both modes (\code{HallCCryoTarget}\simcidx{Classes and interfaces!HallCCryoTarget@\code{HallCCryoTarget}}, members win\_*) \\
130 & D-062 G5 drift\_\allowbreak{}to\_\allowbreak{}cal & dbase.\allowbreak{}f:\allowbreak{}1104 & \code{drift\_\allowbreak{}to\_\allowbreak{}cal}\simcidx{Parameters!drift_to_cal@\code{drift\_to\_cal}} defaults to 0 and is not checked: arms 7/8 without it put the calorimeter face at the target plane & Fortran kept & Fortran kept & calorimeter runs without drift\_\allowbreak{}to\_\allowbreak{}cal & Question: refuse 0? kept (\code{calo} plug-in) \\
131 & D-062 G8 π0 decay in tail 1 & event.\allowbreak{}f:\allowbreak{}899-901 (complete\_\allowbreak{}ev),\allowbreak{} radc.f tail 1 & pizero\_\allowbreak{}decay is called again when the radiative tail 1 re-calls complete\_\allowbreak{}ev: a second, different photon pair per event (as D-038 F9) & Fortran kept & Fortran kept & π0 runs with radiation (photon pair of \textasciitilde{}5 \% of the events) & Question: keep the first pair? kept \\
132 & D-065 P1 π0 in HMS/SHMS & simc.\allowbreak{}f:\allowbreak{}1403-1609,\allowbreak{} mc\_\allowbreak{}hms/\allowbreak{}mc\_\allowbreak{}shms & with doing\_\allowbreak{}pizero and hadron\_arm 1 or 5 the π0 itself (m2 = Mpi0²) is transported through the magnetic spectrometer like a charged particle (apertures, COSY maps, hut, multiple scattering); pizero\_\allowbreak{}decay still fills the photon columns, which no spectrometer uses; RECON\%p = SP\%p (D-063). With doing\_\allowbreak{}decay = 1 the π0 reaches the project/\allowbreak{}transp decay test and STOPs (as D-062 G1) & Fortran kept & Fortran kept & π0 runs with hadron\_arm 1/5 only (reference/\allowbreak{}runs/\allowbreak{}p15f1\_\allowbreak{}pi0\_\allowbreak{}\{shms,\allowbreak{}hms\}\_\allowbreak{}A); \code{cli\_\allowbreak{}shipped\_\allowbreak{}inputs} & Question: is "π0 kinematics in the HMS/SHMS acceptance" a meaningful mode, or should it be refused / use a neutral-particle acceptance? Kept in both modes; the corrected mode never decays the π0 in flight (doing\_\allowbreak{}decay = 1 runs, D-037 F9), -{}-fortran-compat STOPs \\
133 & D-067 T4 field step limit & trg\_\allowbreak{}track.\allowbreak{}f:\allowbreak{}210 & \code{max\_\allowbreak{}steps\ =\ int(\allowbreak{}max(\allowbreak{}dist0,\allowbreak{}10*dl)/\allowbreak{}dl)*10} with the SIGNED dist0: arms called with −sinθ (SOS, HRS-L, SHMS) always get 100 iterations; in the shipped example 208 of 8174 (2.5 \%) SHMS events that passed mc\_shms are rejected by this limit (the iterated reconstruction had run away) & Fortran kept & Fortran kept & using\_\allowbreak{}tgt\_\allowbreak{}field,\allowbreak{} left-side arms & Question: abs(dist0)? kept in both modes \\
134 & D-067 T5 field convergence & trg\_\allowbreak{}track.\allowbreak{}f:\allowbreak{}794-865 & track\_\allowbreak{}to\_\allowbreak{}tgt tests delx of the START of each iteration: the final result is never tested; after 10 iterations it is accepted if that delx ≤ 0.2 (134 captured calls) & Fortran kept & Fortran kept & using\_\allowbreak{}tgt\_\allowbreak{}field reconstruction & Question: test the final delx? kept \\
135 & D-067 T6 field factor & trg\_\allowbreak{}track.\allowbreak{}f:\allowbreak{}81,\allowbreak{} 123, 192 & factor = 90/E instead of c²·1e-16/\allowbreak{}1e6 = 89.8755178737: bending 0.14 \% too strong & Fortran kept & Fortran kept & every using\_\allowbreak{}tgt\_\allowbreak{}field run & Question: use 89.8755? kept \\
136 & D-067 T13 π0 bent in the field & dbase.\allowbreak{}f:\allowbreak{}330-346 & sign\_\allowbreak{}hadron = ±1 for doing\_\allowbreak{}pizero: the neutral π0 is bent in the target field as a charged particle (adds to D-065 P1) & Fortran kept & Fortran kept & π0 + field & Question: kept \\
137 & D-067 T15 polarization vector & event.\allowbreak{}f:\allowbreak{}811-813,\allowbreak{} 851-853, 1206-1208, 1248-1250 & the polarization for the event angles is (−pol·sin|B|,\allowbreak{} 0, pol·cos|B|): always beam left and in-plane, targ\_Bphi and the sign of targ\_\allowbreak{}Bangle ignored (in the shipped example the angle frame and the tracking field are on opposite sides) & Fortran kept & Fortran kept & the 8 polarized-angle columns & Question: kept \\
138 & D-060/D-061 F2 SOS 7.62 cm & mc\_\allowbreak{}sos.\allowbreak{}f:\allowbreak{}337-340,\allowbreak{} 349-360 & the last 7.62 cm to the new exit window are a straight line without \code{project}: no pathlength (FP path −7.62 cm), no decay check over that segment; the code comment says "WE REALLY NEED TO FIX THIS UP" & Fortran kept & Fortran kept & SOS FP path (TOF), decay ON runs & Question: include the 7.62 cm in pathlen / decay? \\
139 & D-061 F3 SOS/HRS recon fry & simc.\allowbreak{}f:\allowbreak{}1349-1355,\allowbreak{} 1472, 1738 & mc\_sos (and mc\_hrs*) receive the raster \code{fry} (−rastery, 0 without correct\_\allowbreak{}raster) as the recon vertical position; mc\_\allowbreak{}hms/\allowbreak{}mc\_\allowbreak{}shms receive xtar\_init (the true x at z = 0) & Fortran kept & Fortran kept & SOS ytar/delta recon with an extended beam spot & Question: intended difference (SOS/HRS matrices fitted with the raster), or a leftover? \\
140 & D-061 F6 HRS aerogel & physics\_\allowbreak{}kaon.\allowbreak{}f:\allowbreak{}155-157,\allowbreak{} semi\_\allowbreak{}physics.\allowbreak{}f:\allowbreak{}617-619 & the survival weight of HRS kaons uses zaero = −183 cm ("aero at 130cm, last project=313(S2)"),\allowbreak{} but mc\_\allowbreak{}hrs*\_\allowbreak{}hut ends the pathlength at the calorimeter back (432.3 cm; 432.3 − 183 = 249.3 ≠ 130) & Fortran kept & Fortran kept & HRS kaon / SIDIS-kaon weights without decay (only visible with D-061 F1 corrected) & Question: which distance is intended? \\
141 & U-A3 pion 1/Z (39d7de0) & upstream 59ccdf6 event.\allowbreak{}f:\allowbreak{}1469-1501 & which\_pion 2/3: the Clebsch-Gordan factors get targ\%\allowbreak{}Z/\allowbreak{}targ\%\allowbreak{}N,\allowbreak{} but only in the hydpi/\allowbreak{}deutpi branches; tgtweight is now set only for which\_pion 0/10 and 1/11, so for A ≥ 3 (doing\_\allowbreak{}hepi) which\_pion 2/3 the weight is per nucleon × 1 (NEW/OLD = 1/Z, measured on C-12: 1/6) & Fortran kept & Fortran kept & pion which\_pion 2/3 on A ≥ 3 only (reference/\allowbreak{}runs/\allowbreak{}up\_\allowbreak{}upf1\_\allowbreak{}pion\_\allowbreak{}c12\_\allowbreak{}wp\{2,\allowbreak{}3\}\_\allowbreak{}A) & Question: intended (deutpi.\allowbreak{}or.\allowbreak{}hepi branch missing)? D-082: the author is asked; until then SIMC++ keeps the 294f662 weight in both modes (\code{-{}-{}fortran-{}compat}\simcidx{Command-line options!fortran-compat@\code{-{}-{}fortran-{}compat}} cannot reproduce 59ccdf6 here) \\
142 & U-B1 eep Wi/epsiloni & upstream 59ccdf6 results\_\allowbreak{}write.\allowbreak{}f (new eep columns), event.\allowbreak{}f:\allowbreak{}723-726 & the new eep columns Wi and epsiloni are always 0 (main\%\allowbreak{}W/\allowbreak{}epsilon are set only in the meson block); for rho Wi/ti are exposed in GeV/1000 (D-056 R5) & Fortran kept & Fortran kept & the upstream-layout eep/rho trees only & Question: intended? Kept in both modes in the upstream layout (the SIMC++ default since U.2; columns 36 \code{Wi}, 38 \code{epsiloni}), proved bitwise by l1\_\allowbreak{}simcrun\_\allowbreak{}ranlux\_\allowbreak{}up\_\allowbreak{}eep\_\allowbreak{}* \\
143 & U-B2 nuclear elastic model & upstream nuclear\_\allowbreak{}elastic.\allowbreak{}f (4df7351); results\_\allowbreak{}write.\allowbreak{}f:\allowbreak{}54 (20d46fc) & sig\_\allowbreak{}nuc\_\allowbreak{}elastic: W1 = 0 (spin-0 nucleus only), Gaussian Fgauss with R = 1.07 A\textasciicircum{}1/\allowbreak{}3 fm except A = 205/\allowbreak{}56/\allowbreak{}28/\allowbreak{}27; in the run W is computed with Mp; corrsing and Pm\_Heep* are never computed for doing\_\allowbreak{}nuc\_\allowbreak{}elast (0); the shipped ee\_\allowbreak{}calcium48 kinematics gives wtcontr = 0 (no event passes pass\_cuts) & Fortran kept & Fortran kept & doing\_\allowbreak{}nuc\_\allowbreak{}elast runs only & Question: is the model meant for spin-0 targets only? Kept in both modes (\code{nuclear::\allowbreak{}Sig\allowbreak{}Nuc\allowbreak{}Elastic}, \code{FillEepNtuple}) \\
\end{simctable}
\end{landscape}
\subsection*{\texorpdfstring{Category C: acceptance mode (open)}{Category C: acceptance mode (open)}}\label{translation_report-06u5fcatalogue:category-c-acceptance-mode-open}
\begin{landscape}
\begin{simctable}{scriptsize}{8}{l>{\raggedright\arraybackslash}p{\dimexpr 0.082\linewidth-2\tabcolsep\relax}>{\raggedright\arraybackslash}p{\dimexpr 0.123\linewidth-2\tabcolsep\relax}>{\raggedright\arraybackslash}p{\dimexpr 0.196\linewidth-2\tabcolsep\relax}>{\raggedright\arraybackslash}p{\dimexpr 0.112\linewidth-2\tabcolsep\relax}>{\raggedright\arraybackslash}p{\dimexpr 0.102\linewidth-2\tabcolsep\relax}>{\raggedright\arraybackslash}p{\dimexpr 0.123\linewidth-2\tabcolsep\relax}>{\raggedright\arraybackslash}p{\dimexpr 0.139\linewidth-2\tabcolsep\relax}}
%
\toprule\rowcolor{simchead}\simcth{\#} & \simcth{id} & \simcth{Fortran location} & \simcth{Fortran behaviour} & \simcth{default (corrected) mode} & \simcth{\code{-{}-{}fortran-{}compat}} & \simcth{effect} & \simcth{question / marker} \\\midrule\endfirsthead
\multicolumn{8}{@{}l}{\simccontinued}\\\toprule\rowcolor{simchead}\simcth{\#} & \simcth{id} & \simcth{Fortran location} & \simcth{Fortran behaviour} & \simcth{default (corrected) mode} & \simcth{\code{-{}-{}fortran-{}compat}} & \simcth{effect} & \simcth{question / marker} \\\midrule\endhead
\midrule\multicolumn{8}{r@{}}{\simccontinues}\\\endfoot
\bottomrule\endlastfoot
144 & D-009 elastic in phsp & dbase.\allowbreak{}f:\allowbreak{}217,\allowbreak{} 227-234; event.f:546 vs 709 & hydrogen + \code{doing\_\allowbreak{}phsp}\simcidx{Parameters!doing_phsp@\code{doing\_phsp}} keeps \code{doing\_\allowbreak{}hyd\_\allowbreak{}elast}, so \code{complete\_\allowbreak{}ev} uses elastic kinematics (not phase space), with sigcc = 1 and normfac = 1 & Fortran kept & Fortran kept & all H acceptance runs & \code{FORTRAN-{}OPEN} in Pipeline.\allowbreak{}cpp \\
145 & D-009 radiation in phsp & dbase.\allowbreak{}f:\allowbreak{}227-228 & only \code{rad\_\allowbreak{}flag}\simcidx{Parameters!rad_flag@\code{rad\_flag}} is zeroed; \code{using\_\allowbreak{}rad}\simcidx{Parameters!using_rad@\code{using\_rad}} stays on: radiative tails are generated in acceptance mode & Fortran kept & Fortran kept & P1.2 phsp Egamma up to 241.9 MeV & \code{FORTRAN-{}OPEN} in SetupBuilder.\allowbreak{}cpp \\
146 & D-028 F1 target mass & dbase.\allowbreak{}f:\allowbreak{}227-233 vs 293-295; init.\allowbreak{}f:\allowbreak{}411-413; radc.f:297 & phsp clears \code{doing\_\allowbreak{}eep} before the masses are set, so \code{Mtar\_\allowbreak{}struck\ =\ 0} and \code{Egamma\_\allowbreak{}tot\_\allowbreak{}max} is one proton mass too low & Fortran kept & Fortran kept & P1.2 phsp Egamma\_max 241.858 instead of 1180.130 MeV; at 6.4 GeV (acc\_\allowbreak{}phsp\_\allowbreak{}K2) negative → 0 events & \code{FORTRAN-{}OPEN} in SetupBuilder.\allowbreak{}cpp \\
\end{simctable}
\end{landscape}
\subsection*{\texorpdfstring{Category D: faithful quirks (kept)}{Category D: faithful quirks (kept)}}\label{translation_report-06u5fcatalogue:category-d-faithful-quirks-kept}
\begin{simctable}{footnotesize}{5}{l>{\raggedright\arraybackslash}p{\dimexpr 0.167\linewidth-2\tabcolsep\relax}>{\raggedright\arraybackslash}p{\dimexpr 0.208\linewidth-2\tabcolsep\relax}>{\raggedright\arraybackslash}p{\dimexpr 0.255\linewidth-2\tabcolsep\relax}>{\raggedright\arraybackslash}p{\dimexpr 0.214\linewidth-2\tabcolsep\relax}}
\caption{Category D: faithful quirks (kept)}\\%
\toprule\rowcolor{simchead}\simcth{\#} & \simcth{id} & \simcth{Fortran location} & \simcth{Fortran behaviour (kept in both modes)} & \simcth{why kept} \\\midrule\endfirsthead
\multicolumn{5}{@{}l}{\simccontinued}\\\toprule\rowcolor{simchead}\simcth{\#} & \simcth{id} & \simcth{Fortran location} & \simcth{Fortran behaviour (kept in both modes)} & \simcth{why kept} \\\midrule\endhead
\midrule\multicolumn{5}{r@{}}{\simccontinues}\\\endfoot
\bottomrule\endlastfoot
147 & D-018 CTP float & CTP/\allowbreak{}thTestParse.\allowbreak{}c:\allowbreak{}392-393 & every real .inp literal is stored as a 32-bit float & defines the run inputs \\
148 & D-010 literals & -fdefault-real-8; physics\_\allowbreak{}proton.\allowbreak{}f & double literals, \code{pow} for real exponents, x+y+z dot products & bitwise numerics \\
149 & D-024/D-018 powers & gfortran & \code{x**2} → x·x; \code{x**n} (integer n) → \code{\_\allowbreak{}\_\allowbreak{}powidf2} (\code{optics::\allowbreak{}powi}) & bitwise numerics \\
150 & D-015 F1 / D-021 lfit & cern/\allowbreak{}lfit.\allowbreak{}f,\allowbreak{} mc\_\allowbreak{}*\_\allowbreak{}hut.\allowbreak{}f & focal plane real\emph{4, accumulators real}8, a hit at exactly 0.0 ignored (KEY = 0) & CERNLIB behaviour \\
151 & D-015 F3 closure & COSY maps & forward→recon closure not exact (HMS xptar RMS 1 mr) & optics \\
152 & D-015 fry & mc\_\allowbreak{}hms\_\allowbreak{}recon.\allowbreak{}f:\allowbreak{}118,\allowbreak{} mc\_\allowbreak{}shms\_\allowbreak{}recon.\allowbreak{}f:\allowbreak{}136 & fry = 0 replaced by 1e-30 & avoids 0**0 \\
153 & D-061 F4 SOS recon guard & mc\_\allowbreak{}sos\_\allowbreak{}recon.\allowbreak{}f:\allowbreak{}120-122 & every hut value (xfp, xpfp, yfp, ypfp, fry) with |v| ≤ 1e-30 is replaced by 1e-30 (HMS/SHMS: only fry) & \code{Cosy\allowbreak{}Engine::\allowbreak{}recon(\allowbreak{}2,\ .\allowbreak{}.\allowbreak{}.\allowbreak{})} \\
154 & D-061 F2 SOS hut start & mc\_\allowbreak{}sos\_\allowbreak{}hut.\allowbreak{}f:\allowbreak{}283-284,\allowbreak{} mc\_\allowbreak{}sos.\allowbreak{}f:\allowbreak{}369 & the hut starts at zinit = −3.206267 cm, after the exit foil at −3.22 cm: the first drift is negative and clamped to 0.001 cm & faithful (\code{McSosHut.\allowbreak{}cpp}) \\
155 & D-061 F7 SOS cal & mc\_\allowbreak{}sos\_\allowbreak{}hut.\allowbreak{}f:\allowbreak{}519-525,\allowbreak{} 535-539 & the calorimeter aperture and fiducial cuts are commented out: sSTOP\_cal is never incremented & faithful \\
156 & D-062 G4 π0 FP & simc.\allowbreak{}f:\allowbreak{}1599-1603 & FP\%p (ntuple *fp of the hadron arm) = the last photon that hit (mc\_calo writes x\_fp only on success); RECON\%p = SP\%p (no π0 reconstruction); xcal/\allowbreak{}ycal\_\allowbreak{}gamma = −1e10 for a missed photon & faithful (\code{driver::\allowbreak{}Pizero\allowbreak{}Calorimeter}\simcidx{Classes and interfaces!PizeroCalorimeter@\code{PizeroCalorimeter}}) \\
157 & D-062 G6/G7 mc\_calo & simc.\allowbreak{}f:\allowbreak{}1756; calo/\allowbreak{}mc\_\allowbreak{}calo.\allowbreak{}f & dead test \code{abs(\allowbreak{}spec\%\allowbreak{}p\%\allowbreak{}phi-{}pi/\allowbreak{}2)\ .\allowbreak{}eq.\ 10.} (zhadron unused); mc\_calo ignores ms\_flag, wcs\_flag, resmult, frx, fry, th\_spec, spect\_side (no smearing, no resolution factor: ntup\%\allowbreak{}resfac stays 0 / the stale tmpfact for the electron arm); mc\_\allowbreak{}calo\_\allowbreak{}recon is never called & faithful; mc\_\allowbreak{}calo\_\allowbreak{}recon ported as a library function for its L1 table \\
158 & D-062 G9/G10 & infiles/\allowbreak{}nps\_\allowbreak{}*.inp; calo/, pizero\_\allowbreak{}decay.\allowbreak{}f & shipped NPS inputs have no random\_\allowbreak{}seed and nps\_sidis sets doing\_pion and doing\_semi (SIDIS wins); SIDIS π0 CENTRAL.sigcc = 0 (z = 1.0003); no REAL(16) argument literals in the calorimeter/\allowbreak{}π0 code & documented \\
159 & D-065 S1 off-arm columns & simc.\allowbreak{}f:\allowbreak{}1612-1618,\allowbreak{} 1806-1812 & with spect\_mode 1/−1/−2 the switched-off arm gets RECON delta/\allowbreak{}yptar/\allowbreak{}xptar = SP, no aperture test, and its FP (\code{*xfp\ *xpfp\ *yfp\ *ypfp}) and recon \code{*ytar} columns stay 0.0 in every event; its .geni counters are 0 & faithful in both modes; docs/\allowbreak{}manual/\allowbreak{}markdown/\allowbreak{}reference/\allowbreak{}OUTPUT\_\allowbreak{}FORMAT.md §1 \\
160 & D-065 S2 spect\_mode values & dbase.\allowbreak{}f:\allowbreak{}531-532,\allowbreak{} 949-950 & any spect\_mode other than 1/−1/−2 means "both arms on", unchecked & faithful; the corrected mode logs a WARNING for values outside \{0, 1, −1, −2\} (SetupBuilder,\allowbreak{} \code{cli\_\allowbreak{}shipped\_\allowbreak{}inputs}) \\
161 & D-065 S3 slop & init.\allowbreak{}f:\allowbreak{}159-200 & the MC slop is added to the generation window only for arms whose MC is on (genvol 0.0195475 vs 0.0224744 for eep\_\allowbreak{}h\_\allowbreak{}elastic\_\allowbreak{}rad\_\allowbreak{}A) & faithful \\
162 & D-065 W1-W3 shipped inputs & simc\_\allowbreak{}gfortran/\allowbreak{}infiles & 18 of the 32 shipped main inputs use legacy \code{.} names: CTP ignores them (registered as dummy ints) and the run STOPs \code{bad\ targ.\allowbreak{}can\ value} (dbase.\allowbreak{}f:\allowbreak{}461); 3 more run with 20 ignored names from a legacy extra file; temp.inp = mc\_\allowbreak{}neg\_\allowbreak{}h\_\allowbreak{}z37\_\allowbreak{}ptbin.\allowbreak{}inp; fpi12\_\allowbreak{}deltapp STOPs (D-038 F1) & same STOP / same run; simc++ additionally lists the ignored names and points to convert\_\allowbreak{}inputfile.\allowbreak{}sh (§16) \\
163 & D-067 T3/T17 field set-up & simc.\allowbreak{}f:\allowbreak{}129-133,\allowbreak{} 143-147; trg\_\allowbreak{}track.\allowbreak{}f:\allowbreak{}286; simc.\allowbreak{}f:\allowbreak{}1583-1585,\allowbreak{} 1787-1789 & the angle error branches print and continue (angle 0; l.146 prints θp as "p-arm phi"); phi treated as 0; calorimeter arms with the field STOP 'Target field reconstruction not set up for your spectrometer' & faithful (FortranStop for arms 7/8) \\
164 & D-067 T7-T12 trg\_track numerics & trg\_\allowbreak{}track.\allowbreak{}f:\allowbreak{}273,\allowbreak{} 409-416, 184, 200-214, 640-648; simc.\allowbreak{}f:\allowbreak{}1431/\allowbreak{}1579 & pi180 = 3.141592653/\allowbreak{}180; field 0 for z' ≥ 100 or r ≥ 100 cm and for non-finite coordinates (INT() = −2147483648); factor not written when ok = .false.; the loop restarts from the wrong-direction step; SIGN(1.,\allowbreak{}−0.) = −1; track\_\allowbreak{}from\_\allowbreak{}tgt's ok ignored and its z = 0 result discarded; Mh to track\_\allowbreak{}from\_\allowbreak{}tgt but sqrt(m2) after a decay to track\_\allowbreak{}to\_\allowbreak{}tgt & faithful, bitwise (optics::\allowbreak{}TargetField) \\
165 & D-067 T16 targ\_pol = 0 & event.\allowbreak{}f:\allowbreak{}895,\allowbreak{} 1289 & theta\_tarq = acos(0/0) = NaN in every event (column th\_tarq); phi\_targ = betai = 0 & faithful \\
166 & D-061 F5 HRS circles & mc\_\allowbreak{}hrs*.f:\allowbreak{}159-180 & the circular apertures at z = 65.686 and 80.436 cm (before the slit box) count as slit\_hor (stop\_\allowbreak{}where 20, 21) & faithful \\
167 & D-061 F5/F7 HRS cal / VDC offset & mc\_\allowbreak{}hrs*\_\allowbreak{}hut.\allowbreak{}f:\allowbreak{}463-470,\allowbreak{} 480-487; mc\_\allowbreak{}hrsl.\allowbreak{}f:\allowbreak{}525,\allowbreak{} mc\_\allowbreak{}hrsr.\allowbreak{}f:\allowbreak{}524 & the calorimeter cuts are commented out (cal counter always 0); y\_fp gets the VDC offset (−0.78 L, −0.48 R) after ys (the recon input) is set, so the recon ignores it & faithful \\
168 & D-061 F5 HRS stop\_where & g\_\allowbreak{}dump\_\allowbreak{}all\_\allowbreak{}events.\allowbreak{}inc,\allowbreak{} mc\_hrs*.f & COMMON /\allowbreak{}g\_\allowbreak{}dump\_\allowbreak{}all\_\allowbreak{}l|r/\allowbreak{} stop\_\allowbreak{}where,\allowbreak{} x\_stop, y\_stop written at each stop, never read by SIMC & mirrored in the kernel (\code{McHrs::\allowbreak{}dump}\simcidx{Classes and interfaces!McHrs@\code{McHrs}}) for the L1 tables \\
169 & D-061 F4 HRS recon guard & mc\_\allowbreak{}hrs*\_\allowbreak{}recon.\allowbreak{}f:\allowbreak{}118 & only hut(5) (fry) is guarded, against the REAL(16) 1.d-30 (binary128 comparison) & \code{Cosy\allowbreak{}Engine::\allowbreak{}recon(\allowbreak{}3|\allowbreak{}4)} \\
170 & P14.2 HRS sieve & mc\_\allowbreak{}hrs*.f:\allowbreak{}122-131 & the mock sieve (compile-time .false.) draws rt, tht and never applies them & faithful (\code{use\_\allowbreak{}sieve} option) \\
171 & D-061 F7 SOS BM01 & mc\_\allowbreak{}sos.\allowbreak{}f:\allowbreak{}234-235 & the BM01 entrance test is asymmetric in y (\code{yt\ >\ w\_\allowbreak{}bm01} or \code{-{}yt\ >\ w\_\allowbreak{}bm01-{}0.\allowbreak{}05*2.\allowbreak{}54}) & faithful (deliberate-looking 0.05 inch) \\
172 & D-019 gamma & radc.\allowbreak{}f:\allowbreak{}100-115 & the polynomial \code{gamma(\allowbreak{})} is dead code (callers bind to libm tgamma); wrong for x ≥ 2 & \code{std::\allowbreak{}tgamma} used, as the binary \\
173 & P5.F1 Q6, Q8-Q10 & radc.f, brem.f & basicrad itail overwrite, schwinger's exp convention (only dhard used), brem/bremos kinematic simplifications,\allowbreak{} spen ≤ 101 terms, one-time warnings, a TAB in a message & no effect / physics model \\
174 & D-012 F2 & enerloss\_\allowbreak{}new.\allowbreak{}f,\allowbreak{} gauss1.f & gauss1 = 0 for u1 = 0.5 exactly; −2 ln 0 replaced by 1e5 & guard; rate \textasciitilde{}0 with the 53-bit RNG \\
175 & D-012 F3 & enerloss\_\allowbreak{}new.\allowbreak{}f & invalid typeflags keep \code{x} (not used by SIMC); Eloss > KE clipped to KE − 1e-7 & guard \\
176 & D-012 F5 & target.f, musc*.f & gauss1 truncations 3.5σ / 99σ / 10σ & explicit choices \\
177 & D-012 F6 & gauss1.f & s = 1 exactly accepted (g = ±0) & negligible \\
178 & D-012 F7 & event.\allowbreak{}f:\allowbreak{}1572-1616 & physics\_\allowbreak{}angles assumes phi0 = π/2 or 3π/2 (warns otherwise) & all spectrometers satisfy it \\
179 & D-012 F1 & event.\allowbreak{}f:\allowbreak{}1641 & dy ≈ 3e-9 instead of 0 for a true dy = 0 (the NaN part is corrected, §2) & rounding \\
180 & D-024 θ = 0 & event.\allowbreak{}f:\allowbreak{}532-538 & elastic θ = 0 can round 1 ulp above Ein (failure); NaN at q = 0 does not STOP & no spectrometer at 0° \\
181 & D-024 Pm & event.\allowbreak{}f:\allowbreak{}546-560 & H elastic: vertex Pm = 0, so PmPer = 0 & definition \\
182 & P4.F1 Q4 & event.f:325 & the hadron physics\_\allowbreak{}angles call uses the previous try's xptar/yptar (overwritten by complete\_\allowbreak{}ev) & no effect \\
183 & D-008 phsp & event.\allowbreak{}f:\allowbreak{}1132-1136,\allowbreak{} 1446-1448; simc.f:397 & phsp: recon hadron P set to the central momentum, sigcc = 1, normfac = 1 & phsp design \\
184 & P4.F1 Q6 & event.\allowbreak{}f:\allowbreak{}1060-1140 & complete\_\allowbreak{}recon\_\allowbreak{}ev: W with Mp, phi\_pq with acos (complete\_\allowbreak{}ev: atan2), recon Ein assigned twice & definitions \\
185 & D-005 .hist limits & init.\allowbreak{}f:\allowbreak{}209-216 & "Input Spectrometer Limits" printed after the slop widening & output definition \\
186 & D-005 tree & make\_\allowbreak{}root\_\allowbreak{}tree.\allowbreak{}f & every branch Float\_t; ≤ 1e7 events & output format \\
187 & D-033 F6 logsf & sf\_\allowbreak{}lookup.\allowbreak{}f:\allowbreak{}184 & the variable is called \code{logsf} but the Em interpolation is linear (no log is taken); SF < 1e-20 (also negative extrapolations,\allowbreak{} and one negative node of benharsf\_\allowbreak{}4.dat) is set to 0 & the coded algorithm \\
188 & D-033 F8 krel & event.f:977 & \code{ntup\%\allowbreak{}krel} is written from the SAVEd local krel, never assigned on the (e,e'p) branches (0) & not an output of (e,e'p) \\
189 & D-038 F11 krel & event.\allowbreak{}f:\allowbreak{}961-977 & pion: \code{ntup\%\allowbreak{}krel} = the SAVEd local krel, set only for D/A quasifree production; 0 for H and coherent production (not an ntuple column of the pion layout) & not an output \\
190 & D-038 F16 STOP status & dbase.f (all STOPs) & every STOP exits with status 0 & SIMC++ exits 3 (\code{FortranStop}\simcidx{Classes and interfaces!FortranStop@\code{FortranStop}}) or 2 (rejected input), Category DEV D-005 \\
191 & D-033 F8/F9 Benhar limits & init.f:335, dbase.\allowbreak{}f:\allowbreak{}592-621 & with use\_\allowbreak{}benhar\_\allowbreak{}sf,\allowbreak{} theory\_\allowbreak{}init is not called: E\_Fermi = 0 (VERTEXedge\%\allowbreak{}Em\%\allowbreak{}min = 0) and VERTEXedge\%\allowbreak{}Pm\%\allowbreak{}max = 790 MeV/c for every file (He-3/He-4 tables go to 976.8); the neutron columns of benharsf\_\allowbreak{}4/\allowbreak{}12/\allowbreak{}56/\allowbreak{}197 normalise like the proton ones & definitions of the input files \\
192 & D-033 F12 D warning & event.\allowbreak{}f:\allowbreak{}588-593 & the "two solutions" warning of D(e,e'p) is printed only for ntried ≤ 5000 (\code{IReaction::\allowbreak{}begin\allowbreak{}Try}\simcidx{Classes and interfaces!IReaction@\code{IReaction}}) & message only \\
193 & D-033 doing\_\allowbreak{}deuterium & init.f:889 & theory\_\allowbreak{}init overwrites doing\_\allowbreak{}deuterium = (nrhoPm = 1 and E\_Fermi < 1) from the file contents & definition \\
194 & P8.2 targetfac & simc.\allowbreak{}f:\allowbreak{}94-95 & \code{3.\allowbreak{}75914d+6} is DOUBLE PRECISION, promoted to real\emph{16 by -fdefault-real-8: targetfac is evaluated in quad precision (objdump: \_\allowbreak{}\_\allowbreak{}divtf3/\allowbreak{}\_\allowbreak{}\_\allowbreak{}multf3) and rounded to real}8; SIMC++ does the same (\code{\_\allowbreak{}\_\allowbreak{}float128}); the double evaluation differs by 1 ulp for the D capture input & bitwise numerics \\
195 & D-015 HMS cal & hms/\allowbreak{}mc\_\allowbreak{}hms\_\allowbreak{}hut.\allowbreak{}f:\allowbreak{}582-588 & calorimeter cut commented out: hSTOP\_cal never incremented & deliberate in the Fortran \\
196 & D-015 SHMS counters & shms/\allowbreak{}mc\_\allowbreak{}shms*.f & 7 counters never incremented; 5 stop reasons unreachable; stop\_id reused / 0 for slit stops & diagnostics \\
197 & P6.F1 F8, F10-F12 & mc\_\allowbreak{}shms(\_\allowbreak{}hut).f,\allowbreak{} transp.f & \code{spec} real\emph{4/real}8 mismatch (unused), unit 1 closed, TOF terms dropped, pathlen to the back of the calorimeter & no effect / definitions \\
198 & P3.F1 & target.\allowbreak{}f:\allowbreak{}62-151 & \code{goto\ (\allowbreak{}10,\allowbreak{}20,\allowbreak{}30)\ narm} falls through to the beam for other narm & unreachable \\
199 & limits\_\allowbreak{}update & event.f limits\_\allowbreak{}update & \code{orig\%\allowbreak{}e\%\allowbreak{}xptar} range updated twice & harmless duplicate \\
200 & .geni histograms & simc.\allowbreak{}f:\allowbreak{}578-586 & geni buffers printed on the H\%gen bin centres & the axes are the same (init.\allowbreak{}f:\allowbreak{}515-570) \\
201 & calculate\_\allowbreak{}central & simc.\allowbreak{}f:\allowbreak{}1268-1281 & \code{main0} is SAVEd: unset fields are 0 & definition \\
202 & montecarlo & simc.\allowbreak{}f:\allowbreak{}1800-1842 & RECON z not written with the arm MC off; recon e E not Coulomb-corrected & definition \\
203 & P1.6 & physics\_\allowbreak{}proton.\allowbreak{}f & real(16) ABI of real*8 functions; fofa for timelike Q²; deForest at θpq = 0 by rounding & no effect / numerics \\
204 & D-038 F13 & physics\_\allowbreak{}pion.\allowbreak{}f:\allowbreak{}467-516 (sig\_\allowbreak{}param04,\allowbreak{} not called) & pole at |t| = 0.02 GeV² (Inf/NaN), \code{alog} on a real*8 argument (= double log with -fdefault-real-8) & the routine's formula; dead code \\
205 & D-038 F14 & physics\_\allowbreak{}pion.\allowbreak{}f:\allowbreak{}659-669 (sigmaid) & nearest (lower) bin in Q², W and cos θ*, no interpolation; only θ rows 1-6 of the 23 stored per (Q², W) are reachable; σ\_T divided by max(0.2, Q²) and by (W−1.132)/\allowbreak{}0.1 above 1.232 GeV & the coded model (MAID2007 tables of P. Bosted) \\
206 & D-038 pi0 decay & pizero\_\allowbreak{}decay.\allowbreak{}f:\allowbreak{}1-94 & 2 grnd per call (φ* then cos θ*); on \code{Mh/\allowbreak{}2\ <\ Mgamma} (impossible,\allowbreak{} Mgamma = 0) the photons are not written & \code{pion::\allowbreak{}pizeroDecay} returns \code{success\ =\ false} and the caller keeps its values \\
207 & P6.4 1.d-27 & hms/\allowbreak{}pion\_\allowbreak{}coll\_\allowbreak{}absorb.\allowbreak{}f:\allowbreak{}76,\allowbreak{} 82 & \code{sigA*1.\allowbreak{}d-{}27}: a DOUBLE PRECISION literal promoted to real*16 by -fdefault-real-8,\allowbreak{} so the product is formed in quad precision (objdump: \_\_multf3) with the quad value of 1e-27; SIMC++ does the same (\code{\_\allowbreak{}\_\allowbreak{}float128}); the double product differs by 1 ulp in 7 of 788 table rows & bitwise numerics \\
208 & D-037 F10 save\_coll & mc\_\allowbreak{}hms.\allowbreak{}f:\allowbreak{}216,\allowbreak{} mc\_\allowbreak{}shms.\allowbreak{}f:\allowbreak{}517 & \code{save\_\allowbreak{}Hcollsuccess/\allowbreak{}save\_\allowbreak{}Scollsuccess\ =\ 1.\allowbreak{}5} on a collimator stop, never reset or read (dead COMMON /\allowbreak{}save\_\allowbreak{}coll/\allowbreak{}); mirrored as \code{Mc\allowbreak{}Hms/\allowbreak{}Mc\allowbreak{}Shms::\allowbreak{}save\allowbreak{}Coll\allowbreak{}Success(\allowbreak{})} for the L1 tables only & no effect \\
209 & D-037 F12 thick\_temp & mc\_\allowbreak{}hms\_\allowbreak{}coll.\allowbreak{}f:\allowbreak{}79-96,\allowbreak{} mc\_\allowbreak{}shms\_\allowbreak{}coll.\allowbreak{}f:\allowbreak{}81-98 & \code{thick\_\allowbreak{}temp} is SAVEd and set only on a hit step (\code{thick\ =\ thick\_\allowbreak{}temp} is read on every step, used only on hit steps); kept as a member & no effect \\
210 & D-037 F6 & shared/\allowbreak{}transp.\allowbreak{}f:\allowbreak{}158 & (already Category A, D-006 kaon decay) now covered bitwise by transp\_\allowbreak{}decay\_\allowbreak{}hms/\allowbreak{}shms (55/103 rows of the Mk case) & — \\
211 & D-037 F13/F14 & gauss1.f, call\_\allowbreak{}ranlux.\allowbreak{}f; CTP & real*16 return ABI of grnd/gauss1 (harness only); CTP float inputs (D-018) & — \\
212 & D-043 F4 binary128 & physics\_\allowbreak{}kaon.\allowbreak{}o & DOUBLE COMPLEX z1a..z8a, ur, ui, dble() and the literal 1.d+06 are kind 16 under -fdefault-real-8: the Saghai amplitude sums run in binary128 (\code{\_\allowbreak{}\_\allowbreak{}float128}, glibc \code{hypotf128} for abs), px(2) = q22/1.d+06 is rounded quad → float; cern/fint.f computes in REAL(8) (implicit typing), not in single precision as CERNLIB & bitwise numerics (L1 eekeek/\allowbreak{}eekeeks \code{\_\allowbreak{}fintabi} bitwise) \\
213 & D-043 F14 krel & event.\allowbreak{}f:\allowbreak{}961-977 & as D-038 F11: ntup\%krel never set on the H/coherent kaon branch & not an output \\
214 & D-043 F16 STOP & dbase.f, physics\_\allowbreak{}kaon.\allowbreak{}f:\allowbreak{}539 & every STOP exits with status 0 & SIMC++ exits 3 (\code{FortranStop}\simcidx{Classes and interfaces!FortranStop@\code{FortranStop}}) \\
215 & P11.1 tfcos & physics\_\allowbreak{}delta.\allowbreak{}f:\allowbreak{}74 & \code{tfcos-{}1.\ .\allowbreak{}lt.\ 1.\allowbreak{}d-{}8} compares in binary128 (1.d-8 is a REAL(16) literal, \_\_lttf2); SIMC++ does the same (\code{\_\allowbreak{}\_\allowbreak{}float128}, both modes) & bitwise numerics \\
216 & D-048 F9 stale fields & event.\allowbreak{}f:\allowbreak{}961-977,\allowbreak{} physics\_\allowbreak{}delta.\allowbreak{}f & ntup\%krel = the SAVEd local krel, never set for H (0); ntup\%\allowbreak{}sigcm2 is never set by peedelta (not columns of the Delta ntuple); tgtweight = 1 on the Delta branch (event.\allowbreak{}f:\allowbreak{}1442) also for D/A & not outputs / the coded weight \\
217 & D-048 F11 STOP & dbase.\allowbreak{}f:\allowbreak{}779-788 & the Delta STOP exits with status 0 & SIMC++ exits 3 (\code{FortranStop}\simcidx{Classes and interfaces!FortranStop@\code{FortranStop}}, compat only; corrected: no STOP) \\
218 & P12.1 Q1 early return & semi\_\allowbreak{}physics.\allowbreak{}f:\allowbreak{}280-284 & below the mmpi2 threshold peepiX returns 0 before main\%\allowbreak{}davejac,\allowbreak{} ntup\%sigcm and the survival probability (decdist) are set: those ntuple columns keep the previous event's values (ntup\%\allowbreak{}xfermi is already written) & weight 0; faithful stale-value semantics (D-024) \\
219 & D-051 F5 dead inputs & semi\_\allowbreak{}physics.\allowbreak{}f:\allowbreak{}142,\allowbreak{} 525; dbase.\allowbreak{}f:\allowbreak{}1096 & the input pt\_\allowbreak{}b\_\allowbreak{}param is read and ignored (\code{b\ =\ pt\_\allowbreak{}b\_\allowbreak{}param} is commented out); the pt² slope is b = 1/(0.120 z² + 0.200) GeV\textsuperscript{−}² & kept as the Fortran (D-051 F5); SIMC++ prints a WARNING when pt\_\allowbreak{}b\_\allowbreak{}param ≠ 0 (\code{sidis::\allowbreak{}ignored\allowbreak{}Input\allowbreak{}Notes}) \\
220 & D-051 F13 CENTRAL.sigcc & simc.\allowbreak{}f:\allowbreak{}1236-1290 & the central peepiX call uses vertex0 with z and p\_T² from complete\_\allowbreak{}recon\_\allowbreak{}ev: meaningful (unlike pion D-038 F5); its survival factor uses main0 FP path 0; the log prints 'multiplicity' = σ\_had of that call & kept (SidisProduction::\allowbreak{}prepareCentral is a no-op); the SIDIS central log lines (simc.\allowbreak{}f:\allowbreak{}1288-1303) are not printed by SIMC++ \\
221 & D-051 F5 doing\_cent & event.\allowbreak{}f:\allowbreak{}1523; semi\_\allowbreak{}physics.\allowbreak{}f:\allowbreak{}157-210 & the central cross section call peepiX(..., .TRUE.) is commented out: the inputs sigc\_flag, sigc\_nbin, sigc\_\allowbreak{}kin\_\allowbreak{}min/\allowbreak{}max/\allowbreak{}ind are read and never used & the branch is ported (\code{sidis::\allowbreak{}peepi\allowbreak{}X(\allowbreak{}.\allowbreak{}.\allowbreak{}.\allowbreak{},\ doing\_\allowbreak{}cent\ =\ true)}) and never called; SIMC++ prints a WARNING when they are set; main\%\allowbreak{}sigcent / ntuple \code{sigcent}, \code{centjac} are always 0 (P12.2) \\
222 & D-051 F5 semi\_\allowbreak{}dilution & event.\allowbreak{}f:\allowbreak{}1524; semi\_\allowbreak{}dilution.\allowbreak{}f & dead code: the call is commented out (\code{ntup\%\allowbreak{}dilu\ =\ 1.\allowbreak{}0}, not an ntuple column) & ported (\code{sidis::\allowbreak{}semi\allowbreak{}Dilution}), not called; unit-tested only \\
223 & P12.1 Q2 sqesub & semi\_\allowbreak{}physics.\allowbreak{}f:\allowbreak{}288-296; F1F2IN21\_\allowbreak{}v1.0.f:\allowbreak{}1309-1401 & the first peepiX call runs SQESUB, which only reads the deuteron smearing tables f1f2tables/\allowbreak{}f1*-cdbonn.\allowbreak{}dat into SAVEd arrays used by the deuteron routines; F1F2IN21 is called with A = 1 & no effect on any output; not ported (the Fortran run needs the files, SIMC++ does not) \\
224 & P12.1 Q3 resmod cut & F1F2IN21\_\allowbreak{}v1.0.f:\allowbreak{}664,\allowbreak{} 953 & \code{if(\allowbreak{}(\allowbreak{}w-{}mp).\allowbreak{}LT.\allowbreak{}wdif(\allowbreak{}1))\ sig\ =\ 0.\allowbreak{}0} with wdif(1) = w−(mp+mpi): never true & no effect (the xpr thresholds do the cut) \\
225 & P12.1 Q4 Ctq5Pdf Q & cteq5/\allowbreak{}Ctq5Pdf.\allowbreak{}f:\allowbreak{}77-80 & Q < Alambda: 'Q out of range' and the CALLER's Q (peepiX Qgev) is set to Alambda for the remaining PDF calls of the event & Q² < 0.051 GeV² only \\
226 & P12.1 Q5 binary128 & fdss/fdss.o & the D literals of fdss.f are REAL(16): (1.D0−X)**4 and **7 are \code{\_\allowbreak{}gfortran\_\allowbreak{}pow\_\allowbreak{}r16\_\allowbreak{}i4} calls (repeated squaring in binary128), the products fFINT·(1−X)\textsuperscript{n}·X\textasciicircum{}p are binary128 rounded once to REAL(8), FAC·(1.D0−D) in fFINT is a binary128 product; Ctq5Pdf.o and semi\_\allowbreak{}physics.\allowbreak{}o have no binary128 arithmetic (only the exact -ff2c return conversions) & bitwise numerics (\code{\_\allowbreak{}\_\allowbreak{}float128} in src/\allowbreak{}physics/\allowbreak{}sidis/\allowbreak{}Fdss.\allowbreak{}cpp) \\
227 & P12.1 Q6 π survival & event.\allowbreak{}f:\allowbreak{}1560; semi\_\allowbreak{}physics.\allowbreak{}f:\allowbreak{}610-627 & with doing\_\allowbreak{}decay = 0 peepiX writes the survival probability to decdist also for pions, but complete\_\allowbreak{}main applies it to the weight for doing\_\allowbreak{}semika only & faithful (the decdist column of decay-OFF π runs is a survival probability,\allowbreak{} not a distance) \\
228 & P12.1 Q7 xbj clamp & semi\_\allowbreak{}physics.\allowbreak{}f:\allowbreak{}255-258 & x\_Bj > 1 prints 'XBj is too large!' and uses x\_Bj = 1 (every such event) & faithful \\
229 & P12.1 Q8 PartonX x ≤ XV(0) & cteq5/\allowbreak{}Ctq5Pdf.\allowbreak{}f:\allowbreak{}148-156,\allowbreak{} 190-194 & the clamp Jx = max(Jx, 0) is only done in the first call with x < Xmin (SAVEd First): later an x ≤ XV(0) = 0 (x = 0, or a NaN x) gives Jx = −1 and POLINT reads XV(−1) and UPD(J0) outside the arrays, i.e. Al and (for JFL = 0, the b̄, with Jq = 0) QL(MXQ) of COMMON /CtqPar1/ & reproduced by a contiguous COMMON storage (\code{sidis::\allowbreak{}Cteq5Pdf}\simcidx{Classes and interfaces!Cteq5Pdf@\code{Cteq5Pdf}}); x = 0 needs Q² = 0, never reached by peepiX; the P12.F1 x = 0 row is bitwise \\
230 & D-057 F7b rho main\%w/t & rho\_\allowbreak{}physics.\allowbreak{}f:\allowbreak{}113,\allowbreak{} 117 & peerho overwrites main\%w with √s in GeV and main\%t with t in GeV² (complete\_\allowbreak{}ev set them in MeV, MeV²) & faithful: the rho ntuple/\allowbreak{}report see the GeV values as in the Fortran \\
231 & D-057 F2b rho\_decay success & rho\_\allowbreak{}decay.\allowbreak{}f:\allowbreak{}1-160 & rho\_decay never sets success = .true.: the incoming success is returned on every non-failure path & faithful (\code{rho::\allowbreak{}rho\allowbreak{}Decay(\allowbreak{}.\allowbreak{}.\allowbreak{}.\allowbreak{},\ success\_\allowbreak{}in,\ .\allowbreak{}.\allowbreak{}.\allowbreak{})}); the call after a failed try is corrected (F2, §2) \\
232 & D-057 F13 rho weight & event.\allowbreak{}f:\allowbreak{}1016-1022; simc.\allowbreak{}f:\allowbreak{}375-391 & no hadron jacobian, genvol = dΩe · 4π · ΔEe (flat CM generation, peerho = gtpr·dσ/\allowbreak{}dΩcm); the truncated Breit-Wigner sampling (±500 MeV: 90.5 \% of the BW) carries no normalisation factor; davejac is computed but not applied & faithful \\
233 & D-057 F14 rho ntuple & results\_\allowbreak{}write.\allowbreak{}f:\allowbreak{}187-229 & 59 columns (OUTPUT\_\allowbreak{}FORMAT.md §1f): \code{t} is the PION t (Mh2 = Mpi2), sigcent/\allowbreak{}xfermi always 0 (zi/pt2i too in 294f662; filled since U.1 as upstream e4dd175, U-D1), Mhadron in MeV, pfermi NaN for H (D-051 F3: 0 in both modes since U.1, fixed upstream) and ±0 for D/³He, jacobian/\allowbreak{}centjac = davejac/\allowbreak{}johnjac,\allowbreak{} phipqi = the rho's φpq; with using\_\allowbreak{}tgt\_\allowbreak{}field \code{mmnuc} is never filled (D-067 T14, corrected by default) & faithful (pfermi: D-051 F3) \\
234 & D-057 F17 rho CENTRAL.sigcc & simc.\allowbreak{}f:\allowbreak{}1236-1290 & CENTRAL.sigcc = peerho for the SHMS central particle treated as a rho of nominal mass with main0\%\allowbreak{}phi\_\allowbreak{}pq = 0 (as D-038 F5) & kept in both modes (\code{Rho\allowbreak{}Production::\allowbreak{}prepare\allowbreak{}Central}\simcidx{Classes and interfaces!RhoProduction@\code{RhoProduction}} does nothing) \\
235 & D-057 F19 rho pass\_cuts & simc.f:237 (D-006) & the electron delta max in the hadron test changes Σ Weight of the rho runs by −24 … −34 \% & corrected by D-006 (§2) \\
236 & D-057 F22 rho\_decay dead code & rho\_\allowbreak{}decay.\allowbreak{}f:\allowbreak{}54,\allowbreak{} 132-136 & dlen/\allowbreak{}ctaurho are computed and unused; the 'going backwards' test is unreachable after the pzprime test & faithful \\
237 & D-057 diag222 & rho\_\allowbreak{}physics.\allowbreak{}f:\allowbreak{}252,\allowbreak{} 262, 293, 390 & the diagnostics (R < 0, sigt < 0, gtpr ≤ 0, 'crazy sigma\_\allowbreak{}rho') go to unit 222 (fort.222; none written in the L2 runs) & kept in the model (\code{Rho\allowbreak{}Model::\allowbreak{}diagnostics(\allowbreak{})}\simcidx{Classes and interfaces!RhoModel@\code{RhoModel}}, printed with 'rho (fort.222):' when verbose); no fort.222 file \\
238 & D-057 literals & rho\_\allowbreak{}physics.\allowbreak{}f:\allowbreak{}280,\allowbreak{} 384 & 1.d+06 and 1.d3 are REAL(16) under -fdefault-real-8: sig219/\allowbreak{}1.d+06 and davesig/\allowbreak{}1.d3 are binary128 divisions & \code{\_\allowbreak{}\_\allowbreak{}float128} divisions (bitwise numerics) \\
239 & U-D1 upstream report/\allowbreak{}quirks & upstream simc.f:894, event.\allowbreak{}f:\allowbreak{}990-996/\allowbreak{}1338-1342,\allowbreak{} init.f:344, benharsf\_\allowbreak{}64.dat,\allowbreak{} infiles/\allowbreak{}ee\_\allowbreak{}calcium48.inp & the .hist has a new line \code{doing\_\allowbreak{}nuc\_\allowbreak{}elast=\ F} for every run; vertex/\allowbreak{}recon z and pt² are computed for every reaction (no RNG; visible in the rho columns zi/pt2i, 0 before); SIDIS z > 1 prints a WARNING per try; the doing\_\allowbreak{}nuc\_\allowbreak{}elast clause of init.f:344 is unreachable (init.f:321 catches it); benharsf\_\allowbreak{}64.dat has identical proton and neutron columns; ee\_\allowbreak{}calcium48.inp has \code{Ebeam\ =\ 2118.\allowbreak{}0.} (parsed as 2118.0), no random\_\allowbreak{}seed and mrec\_amu = mass\_amu (replaced by M − Mp at start-up) & followed by SIMC++ in both modes since U.1 (z/pt²: \code{EepReactionBase}\simcidx{Classes and interfaces!EepReactionBase@\code{EepReactionBase}}, \code{PionProduction}\simcidx{Classes and interfaces!PionProduction@\code{PionProduction}}, \code{RhoProduction}\simcidx{Classes and interfaces!RhoProduction@\code{RhoProduction}}; line: \code{RunReportHist}) \\
\end{simctable}
\subsection*{\texorpdfstring{Category DEV: deliberate design deviations}{Category DEV: deliberate design deviations}}\label{translation_report-06u5fcatalogue:category-dev-deliberate-design-deviations}
\begin{simctable}{footnotesize}{4}{l>{\raggedright\arraybackslash}p{\dimexpr 0.126\linewidth-2\tabcolsep\relax}>{\raggedright\arraybackslash}p{\dimexpr 0.317\linewidth-2\tabcolsep\relax}>{\raggedright\arraybackslash}p{\dimexpr 0.420\linewidth-2\tabcolsep\relax}}
\caption{Category DEV: deliberate design deviations}\\%
\toprule\rowcolor{simchead}\simcth{\#} & \simcth{id} & \simcth{Fortran} & \simcth{SIMC++ (both modes)} \\\midrule\endfirsthead
\multicolumn{4}{@{}l}{\simccontinued}\\\toprule\rowcolor{simchead}\simcth{\#} & \simcth{id} & \simcth{Fortran} & \simcth{SIMC++ (both modes)} \\\midrule\endhead
\midrule\multicolumn{4}{r@{}}{\simccontinues}\\\endfoot
\bottomrule\endlastfoot
240 & D-002 RNG & RANLUX (call\_\allowbreak{}ranlux.\allowbreak{}f; skips every 1000th number) & \code{mt64} (std::\allowbreak{}mt19937\_\allowbreak{}64); \code{random\_\allowbreak{}state\_\allowbreak{}file}\simcidx{Parameters!random_state_file@\code{random\_state\_file}} ignored; validation is statistical. Bitwise replays use \code{ScriptedRandom}\simcidx{Classes and interfaces!ScriptedRandom@\code{ScriptedRandom}} with captured streams. \\
241 & D-005 STOP & \code{STOP} exits with status 0 & \code{simc::\allowbreak{}FortranStop}\simcidx{Classes and interfaces!FortranStop@\code{FortranStop}}, exit status 3; rejected inputs exit 2 \\
242 & D-003 input & CTP .inp, \code{\%} names only (obsolete \code{.} names silently ignored) & .inp (same CTP semantics, unknown names reported like CTP) or JSON; \code{inp2json} \\
243 & D-032 & arm codes as raw integers, 6 half-supported & one \code{SpectrometerId} table; 6 rejected (see §2) \\
244 & not ported & — & nothing: the SOS (arm 2, P14.1), the HRS-R/L (arms 3/4, P14.2), the calorimeter arms 7/8 with π0 → 2γ (P14.3) and the polarized-target field (using\_\allowbreak{}tgt\_\allowbreak{}field,\allowbreak{} P15.2b) are ported \\
245 & D-057 F16 rho A ≥ 4 & dbase.\allowbreak{}f:\allowbreak{}209-214,\allowbreak{} 875-876: no doing\_*rho flag for nint(A) ≥ 4, the run STOPs (exit 0) after printing E\_bind = −1e10; hadron\_arm 6/7/8 STOP in rho\_decay & refused at input in both modes (std::\allowbreak{}invalid\_\allowbreak{}argument; A ≥ 4: SetupBuilder; arms 7/8: the "rho" model; 6 is rejected for every run, D-032); P15.2: the A-range refusal comes after the dbase.\allowbreak{}f:\allowbreak{}461/\allowbreak{}465 target STOPs, so rhotest.inp (legacy names, A = 0) gives the Fortran \code{bad\ targ.\allowbreak{}can\ value} \\
246 & P13.1 rho + other reaction & dbase.\allowbreak{}f:\allowbreak{}136-215 takes the first branch of its if-chain (pion, kaon, delta, semi) and leaves doing\_rho set, so event.\allowbreak{}f:\allowbreak{}420-421 decays that hadron as a rho & doing\_rho together with doing\_\allowbreak{}pion/\allowbreak{}kaon/\allowbreak{}delta/\allowbreak{}semi is refused at start-up (std::\allowbreak{}invalid\_\allowbreak{}argument,\allowbreak{} both modes); doing\_rho with doing\_\allowbreak{}pizero too (results\_\allowbreak{}write.\allowbreak{}f:\allowbreak{}230-236 overwrites the rho columns) \\
247 & P13.2 RANLUX & grnd = RANLUX luxury 3 (cern/\allowbreak{}ranlux.\allowbreak{}f,\allowbreak{} call\_\allowbreak{}ranlux.\allowbreak{}f: 999 of every block of 1000), gauss1.f & besides the default \code{mt64} (D-002), the IRandom \code{"rng":\ \{"type":\ "ranlux"\}} reproduces the Fortran stream bit for bit from random\_\allowbreak{}seed (or \code{"state"}: the 25 integers of *\_\allowbreak{}start\_\allowbreak{}random\_\allowbreak{}state.\allowbreak{}dat): with \code{-{}-{}fortran-{}compat}\simcidx{Command-line options!fortran-compat@\code{-{}-{}fortran-{}compat}} a SIMC++ run can then be identical to a Fortran run (both P13.F1 rho captures: every try, tree, .hist/\allowbreak{}.geni/\allowbreak{}.gen) \\
248 & P12.1 SIDIS files & Ctq5Pdf.\allowbreak{}f:\allowbreak{}SetCtq5 opens cteq5/\allowbreak{}cteq5m.\allowbreak{}tbl,\allowbreak{} fdss.f opens fdss/\allowbreak{}KANLO.GRID (kaons) from the working directory at the first peepiX call & the "sidis" model reads the copies in \code{data/\allowbreak{}sidis} (md5-checked,\allowbreak{} \code{l1\_\allowbreak{}sidis\_\allowbreak{}data\_\allowbreak{}files}; JSON \code{physics.\allowbreak{}data\_\allowbreak{}dir}) at start-up; a missing file is an error at start-up \\
249 & D-033 F1 & the spectral-function file follows from nint(A) only & also selectable in the JSON (\code{pipeline.\allowbreak{}reaction.\allowbreak{}spectral\_\allowbreak{}function}: type and file); a warning when the Fortran C-12 fallback is used \\
250 & outputs & .hist/\allowbreak{}.geni/\allowbreak{}.gen,\allowbreak{} .bin, .root & the same text files, \code{.\allowbreak{}root} (h10, the branches of results\_\allowbreak{}write.\allowbreak{}f for the reaction, Float\_t as make\_\allowbreak{}root\_\allowbreak{}tree),\allowbreak{} plus \code{.\allowbreak{}summary.\allowbreak{}json/\allowbreak{}.\allowbreak{}txt} and \code{.\allowbreak{}gen.\allowbreak{}root}; no \code{.\allowbreak{}bin} \\
251 & P10.1 kaon files & dbase.\allowbreak{}f:\allowbreak{}624-679 reads weightc.\allowbreak{}dat,\allowbreak{} weightd.\allowbreak{}dat,\allowbreak{} saghai\_\allowbreak{}proton.\allowbreak{}dat,\allowbreak{} saghai\_\allowbreak{}sigma0.dat from the working directory & the "kaon" model reads the copies in \code{data/\allowbreak{}kaon} (md5-checked,\allowbreak{} \code{l1\_\allowbreak{}kaon\_\allowbreak{}data\_\allowbreak{}files}; JSON \code{physics.\allowbreak{}data\_\allowbreak{}dir}) at start-up \\
252 & P9.1 MAID files & sigmaid reads maid*.dat from the working directory at its first call (W < 2 GeV), sig\_\allowbreak{}multipole MULTIPOL.file & the "pion" model reads the copies in \code{data/\allowbreak{}pion} (md5-checked,\allowbreak{} \code{l1\_\allowbreak{}pion\_\allowbreak{}data\_\allowbreak{}files}; JSON \code{physics.\allowbreak{}data\_\allowbreak{}dir}) when it is initialised: a missing file is an error at start-up, not in the middle of a run \\
253 & P9.1 peepi writes main & peepi(vertex,\allowbreak{} main) writes main\%\allowbreak{}thetacm,\allowbreak{} phicm, pcm, davejac, johnjac, wcm and ntup\%\allowbreak{}sigcm* & \code{IPhysics\allowbreak{}Model::\allowbreak{}sigcc\allowbreak{}Main(\allowbreak{}vertex,\ main\&)} does the same; \code{sigcc(\allowbreak{}vertex,\ const\ main\&)} evaluates on a copy. COMMON /pfermi/ is \code{MainEvent::\allowbreak{}pfermi}\simcidx{Classes and interfaces!MainEvent@\code{MainEvent}}, ntup\%\allowbreak{}sigcm1/\allowbreak{}2 are \code{Ntup\allowbreak{}Vars::\allowbreak{}sigcm1/\allowbreak{}2} \\
254 & U-DEV1 doing\_\allowbreak{}nuc\_\allowbreak{}elast + another reaction & upstream dbase.\allowbreak{}f:\allowbreak{}136-222 keeps doing\_\allowbreak{}nuc\_\allowbreak{}elast set when doing\_\allowbreak{}pion/\allowbreak{}kaon/\allowbreak{}delta/\allowbreak{}semi/\allowbreak{}rho win the if-chain (and with doing\_\allowbreak{}phsp),\allowbreak{} and event.f then mixes the elastic kinematics into that reaction & such an input is refused at set-up (\code{Derive\allowbreak{}Reaction\allowbreak{}Flags}: "doing\_\allowbreak{}nuc\_\allowbreak{}elast cannot be combined ..."), both modes; test \code{fortran\_\allowbreak{}diff\_\allowbreak{}U\_\allowbreak{}DEV1\_\allowbreak{}nuc\_\allowbreak{}elast\_\allowbreak{}combination} \\
\end{simctable}
\section{\texorpdfstring{The expert list}{The expert list}}\label{translation_report-06u5fcatalogue:63-the-expert-list}
The Category B items are the open physics questions for the user and an expert (P15.X). Each keeps the Fortran behaviour in both modes. Together with the acceptance mode (Category C, D-031 flag 1) and the Delta weight (D-048 F4: the doing\_\allowbreak{}delta weight contains no cross section) they are the questions the translation could not answer.

\begin{simctable}{small}{3}{>{\raggedright\arraybackslash}p{\dimexpr 0.266\linewidth-2\tabcolsep\relax}>{\raggedright\arraybackslash}p{\dimexpr 0.280\linewidth-2\tabcolsep\relax}>{\raggedright\arraybackslash}p{\dimexpr 0.386\linewidth-2\tabcolsep\relax}}
\caption{The expert list}\\%
\toprule\rowcolor{simchead}\simcth{id} & \simcth{Fortran location} & \simcth{question} \\\midrule\endfirsthead
\multicolumn{3}{@{}l}{\simccontinued}\\\toprule\rowcolor{simchead}\simcth{id} & \simcth{Fortran location} & \simcth{question} \\\midrule\endhead
\midrule\multicolumn{3}{r@{}}{\simccontinues}\\\endfoot
\bottomrule\endlastfoot
D-033 F4 D jacobian & event.\allowbreak{}f:\allowbreak{}1020-1024 & is the missing jacobian for D intended? \\
D-033 F1 SF fallback & init.\allowbreak{}f:\allowbreak{}839-851,\allowbreak{} dbase.\allowbreak{}f:\allowbreak{}598-616 & which spectral function for He-3/He-4 (theory files exist only for 2/\allowbreak{}12/\allowbreak{}56/\allowbreak{}197)? \\
D-015 F7 HMS hodo 2 & hms/\allowbreak{}mc\_\allowbreak{}hms\_\allowbreak{}hut.\allowbreak{}f:\allowbreak{}537-551 & copy-paste or the same active area on purpose? \\
D-015 F7 HMS Cherenkov & hms/\allowbreak{}mc\_\allowbreak{}hms\_\allowbreak{}hut.\allowbreak{}f:\allowbreak{}523-529 & double counting or a second foil? \\
D-015 F5 SHMS drifts & shms/\allowbreak{}mc\_\allowbreak{}shms.\allowbreak{}f (zd\_d1men 39.47, zd\_d1out 60.68, zd\_fp 307.95) & use the COSY values? \\
D-017 can radius & target.\allowbreak{}f:\allowbreak{}368-376,\allowbreak{} 456-464 vs 108, 228 & which radius? \\
D-012 F4 th\_can & target.\allowbreak{}f:\allowbreak{}236-244 & reject, or clamp? \\
D-019 one\_tail & radc.\allowbreak{}f:\allowbreak{}198-216,\allowbreak{} 234, 358 & restrict \code{frac} to the enabled tails? \\
D-019 doing\_\allowbreak{}proton & init.f:693 & intended? \\
D-019 intcor\_\allowbreak{}mode 0 & init.\allowbreak{}f:\allowbreak{}701-715 & derivative term intended? \\
D-019 rad\_weight & radc.\allowbreak{}f:\allowbreak{}330-333,\allowbreak{} 403-406, 454-457 & intended (peaked weight contains basicrad)? \\
D-019 tail-1 COMMON & radc.\allowbreak{}f:\allowbreak{}322-328,\allowbreak{} 512 & post-radiation hardcorfac intended? \\
D-019 Q4 phi\_ext & radc.\allowbreak{}f:\allowbreak{}566-642 & include phi\_ext? \\
D-019 Q4 dsoft\_int & radc.\allowbreak{}f:\allowbreak{}590-620 & compute them? \\
D-019 Egamma\_\allowbreak{}min(1) & radc.\allowbreak{}f:\allowbreak{}283-286 & set it? \\
D-005 lambda input & init.\allowbreak{}f:\allowbreak{}692-695 & drop the input or honour it? \\
D-024 stale orig & event.\allowbreak{}f:\allowbreak{}400-415 & — (try rejected) \\
D-024 grnd dEbeam=0 & event.f:204 & — \\
D-024 Eloss gaussians & event.\allowbreak{}f:\allowbreak{}193-205,\allowbreak{} trip\_\allowbreak{}thru\_\allowbreak{}target & — \\
D-015 F2 HMS resmult & hms/\allowbreak{}mc\_\allowbreak{}hms\_\allowbreak{}hut.\allowbreak{}f:\allowbreak{}292 & — \\
D-015 tmpfact & simc.f:1334 & intended? \\
D-015 F6 SHMS counters & shms/\allowbreak{}mc\_\allowbreak{}shms\_\allowbreak{}hut.\allowbreak{}f & separate counters? \\
D-011 recon xptar & simc.f:274, 277 vs 256, 262 & which convention? \\
D-015 F4 fry & simc.\allowbreak{}f:\allowbreak{}1467,\allowbreak{} 1487, 1733, 1753 vs 1472-1482 & intended? \\
D-024 H elastic W & event.f:723 block (pion/kaon only) & intended? \\
D-024 ngen<0 & simc.\allowbreak{}f:\allowbreak{}169-199,\allowbreak{} 341-350 & intended? \\
D-018 herho & dbase.f:213 & intended? \\
D-005 sigc\_nbin & semi\_\allowbreak{}physics.\allowbreak{}f:\allowbreak{}173 & required input? \\
D-038 F4 coherent model & physics\_\allowbreak{}pion.\allowbreak{}f:\allowbreak{}807-854 (exclfit), event.\allowbreak{}f:\allowbreak{}1496-1500 & is a nucleon model meant for coherent production, and should negative cross sections be clamped? \\
D-038 F3 coherent Mrec & dbase.\allowbreak{}f:\allowbreak{}283-288,\allowbreak{} 380-388 & should the input final-nucleus mass be honoured? \\
D-038 F9 Δ re-generation & event.\allowbreak{}f:\allowbreak{}620-631,\allowbreak{} radc.f:324 & intended re-generation? \\
D-038 F12 sig\_\allowbreak{}multipole & physics\_\allowbreak{}pion.\allowbreak{}f:\allowbreak{}197-398 (dead code: the call in peepi is commented out, l.157-163) & single precision and the previous-call amplitudes intended? \\
D-037 F3 coll scattering & hms/\allowbreak{}mc\_\allowbreak{}hms\_\allowbreak{}coll.\allowbreak{}f:\allowbreak{}115-118,\allowbreak{} shms/\allowbreak{}mc\_\allowbreak{}shms\_\allowbreak{}coll.\allowbreak{}f:\allowbreak{}117-120 & is the collimator scattering intended physics (the purpose of the stepping) or should it follow ms\_flag? \\
D-037 F5 pion\_\allowbreak{}coll\_\allowbreak{}absorb & hms/\allowbreak{}pion\_\allowbreak{}coll\_\allowbreak{}absorb.\allowbreak{}f:\allowbreak{}16,\allowbreak{} 30-56 & intended parameters? extrapolate below 85 MeV? \\
D-037 F8 transp decay & shared/\allowbreak{}transp.\allowbreak{}f:\allowbreak{}142-185,\allowbreak{} 232-276 & an intended approximation (the COSY map cannot be split)? \\
D-043 F2 data & saghai\_\allowbreak{}sigma0.dat (content) & which amplitudes belong to s = 2.851, Q² = 0? \\
D-043 F7 decay model & shared/\allowbreak{}transp.\allowbreak{}f:\allowbreak{}149-160,\allowbreak{} 239-251; project.\allowbreak{}f:\allowbreak{}76-85 & intended approximation? \\
D-043 F8 Saghai angles & physics\_\allowbreak{}kaon.\allowbreak{}f:\allowbreak{}19,\allowbreak{} 79-85 & intended (comment "ALWAYS CALCULATING FOR PHI=0")? \\
D-043 F9 coherent kaon & physics\_\allowbreak{}kaon.\allowbreak{}f:\allowbreak{}82-86,\allowbreak{} 215; event.\allowbreak{}f:\allowbreak{}1505 & is a nucleon model meant for coherent production? \\
D-043 F10 Mp & physics\_\allowbreak{}kaon.\allowbreak{}f:\allowbreak{}194-197 & intended? \\
\textbf{D-048 F4 Delta weight} & physics\_\allowbreak{}delta.\allowbreak{}f:\allowbreak{}130-132,\allowbreak{} event.\allowbreak{}f:\allowbreak{}1512 & should SIMC++ offer a real Delta (e,e'p)π cross section, and which one? (no choice can be derived from the Fortran) \\
D-048 F5 two roots & event.\allowbreak{}f:\allowbreak{}682-684 & is a factor 2 missing? \\
D-048 F5b root re-choice & radc.f:322 with event.\allowbreak{}f:\allowbreak{}682-684 & intended re-generation? \\
D-048 F7 recoil mass & dbase.\allowbreak{}f:\allowbreak{}297-300 & π0 mass intended? \\
D-048 F8 which\_pion & event.\allowbreak{}f:\allowbreak{}470-491,\allowbreak{} 612-623; physics\_\allowbreak{}delta.\allowbreak{}f:\allowbreak{}99 & should which\_pion be ignored for doing\_\allowbreak{}delta? \\
D-051 F7 threshold & semi\_\allowbreak{}physics.\allowbreak{}f:\allowbreak{}278-284 & which missing-mass formula was intended? \\
D-051 F8 semi thresholds & event.\allowbreak{}f:\allowbreak{}979-996 & which threshold is intended for kaons / with Fermi motion? \\
D-051 F16 hadron range & init.\allowbreak{}f:\allowbreak{}504-506 & intended generation range? \\
D-051 F10 model choices (R\_h for D) & semi\_\allowbreak{}physics.\allowbreak{}f:\allowbreak{}659-677 & intended for D? (D-051 F10 also lists the slacemcfit and Mtar\_\allowbreak{}struck = Mp choices) \\
D-043 F12 warning & physics\_\allowbreak{}kaon.\allowbreak{}f:\allowbreak{}201-206 & message only \\
D-057 F4 rho nuclei & dbase.\allowbreak{}f:\allowbreak{}563,\allowbreak{} 594; event.\allowbreak{}f:\allowbreak{}336-372; sf\_\allowbreak{}lookup.\allowbreak{}f:\allowbreak{}195-245 & should D/³He rho include Fermi motion / a spectral function? (a model limitation, kept in both modes; SIMC++ computes the ³He Em = 0 without the out-of-bounds read) \\
D-057 F5 rho s & generate\_\allowbreak{}rho.\allowbreak{}f:\allowbreak{}72 & intended? (latent: pfer = 0 always, F4) \\
D-057 F7 rho t & rho\_\allowbreak{}physics.\allowbreak{}f:\allowbreak{}115,\allowbreak{} 233 & model defined at the pole mass on purpose? \\
D-057 F8 rho decay frame & rho\_\allowbreak{}decay.\allowbreak{}f:\allowbreak{}57-77 & intended approximation? \\
D-057 F10 rho xptar & rho\_\allowbreak{}decay.\allowbreak{}f:\allowbreak{}110,\allowbreak{} 142 & approximation accepted? \\
D-057 F11 rho tail 3 & radc.\allowbreak{}f:\allowbreak{}386-387; init.\allowbreak{}f:\allowbreak{}231-232 & should the rho radiate, and with which limits? \\
D-057 F12 rho tail 1 & radc.f:324 & intended? \\
D-012 F4 calo windows & target.\allowbreak{}f:\allowbreak{}73-102,\allowbreak{} 189-222 & which windows / target path should the calorimeter arm get (NPS geometry)? kept in both modes (\code{HallCCryoTarget}\simcidx{Classes and interfaces!HallCCryoTarget@\code{HallCCryoTarget}}, members win\_*) \\
D-062 G5 drift\_\allowbreak{}to\_\allowbreak{}cal & dbase.\allowbreak{}f:\allowbreak{}1104 & refuse 0? kept (\code{calo} plug-in) \\
D-062 G8 π0 decay in tail 1 & event.\allowbreak{}f:\allowbreak{}899-901 (complete\_\allowbreak{}ev),\allowbreak{} radc.f tail 1 & keep the first pair? kept \\
D-065 P1 π0 in HMS/SHMS & simc.\allowbreak{}f:\allowbreak{}1403-1609,\allowbreak{} mc\_\allowbreak{}hms/\allowbreak{}mc\_\allowbreak{}shms & is "π0 kinematics in the HMS/SHMS acceptance" a meaningful mode, or should it be refused / use a neutral-particle acceptance? Kept in both modes; the corrected mode never decays the π0 in flight (doing\_\allowbreak{}decay = 1 runs, D-037 F9), -{}-fortran-compat STOPs \\
D-067 T4 field step limit & trg\_\allowbreak{}track.\allowbreak{}f:\allowbreak{}210 & abs(dist0)? kept in both modes \\
D-067 T5 field convergence & trg\_\allowbreak{}track.\allowbreak{}f:\allowbreak{}794-865 & test the final delx? kept \\
D-067 T6 field factor & trg\_\allowbreak{}track.\allowbreak{}f:\allowbreak{}81,\allowbreak{} 123, 192 & use 89.8755? kept \\
D-067 T13 π0 bent in the field & dbase.\allowbreak{}f:\allowbreak{}330-346 & kept \\
D-067 T15 polarization vector & event.\allowbreak{}f:\allowbreak{}811-813,\allowbreak{} 851-853, 1206-1208, 1248-1250 & kept \\
D-060/D-061 F2 SOS 7.62 cm & mc\_\allowbreak{}sos.\allowbreak{}f:\allowbreak{}337-340,\allowbreak{} 349-360 & include the 7.62 cm in pathlen / decay? \\
D-061 F3 SOS/HRS recon fry & simc.\allowbreak{}f:\allowbreak{}1349-1355,\allowbreak{} 1472, 1738 & intended difference (SOS/HRS matrices fitted with the raster), or a leftover? \\
D-061 F6 HRS aerogel & physics\_\allowbreak{}kaon.\allowbreak{}f:\allowbreak{}155-157,\allowbreak{} semi\_\allowbreak{}physics.\allowbreak{}f:\allowbreak{}617-619 & which distance is intended? \\
U-A3 pion 1/Z (39d7de0) & upstream 59ccdf6 event.\allowbreak{}f:\allowbreak{}1469-1501 & intended (deutpi.\allowbreak{}or.\allowbreak{}hepi branch missing)? D-082: the author is asked; until then SIMC++ keeps the 294f662 weight in both modes (\code{-{}-{}fortran-{}compat}\simcidx{Command-line options!fortran-compat@\code{-{}-{}fortran-{}compat}} cannot reproduce 59ccdf6 here) \\
U-B1 eep Wi/epsiloni & upstream 59ccdf6 results\_\allowbreak{}write.\allowbreak{}f (new eep columns), event.\allowbreak{}f:\allowbreak{}723-726 & intended? Kept in both modes in the upstream layout (the SIMC++ default since U.2; columns 36 \code{Wi}, 38 \code{epsiloni}), proved bitwise by l1\_\allowbreak{}simcrun\_\allowbreak{}ranlux\_\allowbreak{}up\_\allowbreak{}eep\_\allowbreak{}* \\
U-B2 nuclear elastic model & upstream nuclear\_\allowbreak{}elastic.\allowbreak{}f (4df7351); results\_\allowbreak{}write.\allowbreak{}f:\allowbreak{}54 (20d46fc) & is the model meant for spin-0 targets only? Kept in both modes (\code{nuclear::\allowbreak{}Sig\allowbreak{}Nuc\allowbreak{}Elastic}, \code{FillEepNtuple}) \\
\end{simctable}
\section{\texorpdfstring{Cross-check with the decisions}{Cross-check with the decisions}}\label{translation_report-06u5fcatalogue:64-cross-check-with-the-decisions}
Every finding id (\code{F}, \code{G}, \code{T}, \code{R}, \code{S} + number) named in a DECISIONS.md entry about findings was searched in FORTRAN\_\allowbreak{}DIFFERENCES.md together with its decision number: 182 ids checked, 42 without an entry of their own.

Ids without an entry of their own (most are documentation-only findings, Category D, described in the decision text, or ids of the provisional numbering that a later decision superseded, e.g. D-056 R-ids → D-057 F-ids; see the decision for each):

\begin{simctable}{small}{2}{>{\raggedright\arraybackslash}p{\dimexpr 0.105\linewidth-2\tabcolsep\relax}>{\raggedright\arraybackslash}p{\dimexpr 0.367\linewidth-2\tabcolsep\relax}}
\caption{Cross-check with the decisions}\\%
\toprule\rowcolor{simchead}\simcth{decision} & \simcth{finding ids} \\\midrule\endfirsthead
\multicolumn{2}{@{}l}{\simccontinued}\\\toprule\rowcolor{simchead}\simcth{decision} & \simcth{finding ids} \\\midrule\endhead
\midrule\multicolumn{2}{r@{}}{\simccontinues}\\\endfoot
\bottomrule\endlastfoot
D-009 & F1 \\
D-019 & F1 \\
D-024 & F1 \\
D-025 & F2 \\
D-033 & F2 \\
D-035 & F3, F4, F5, F6, F7 \\
D-038 & F6, F7, F8 \\
D-039 & F1, F2 \\
D-043 & F15 \\
D-051 & F11, F17 \\
D-056 & R1, R2, R3, R4, R6, R7, R8, R9 \\
D-057 & R1, R2, R3, R4, R6, R8, R9 \\
D-060 & F1 \\
D-061 & F9, F10 \\
D-062 & F2 \\
D-067 & F2, T8, T9, T10, T11 \\
\end{simctable}
\section{\texorpdfstring{Findings fixed upstream}{Findings fixed upstream}}\label{translation_report-06u5fcatalogue:65-findings-fixed-upstream}
Findings that the newer JLab SIMC (JeffersonLab/\allowbreak{}simc\_\allowbreak{}gfortran,\allowbreak{} followed commit in the file UPSTREAM\_\allowbreak{}SIMC) has fixed itself (D-081). They keep their Category; \code{-{}-{}fortran-{}compat}\simcidx{Command-line options!fortran-compat@\code{-{}-{}fortran-{}compat}} reproduces the fixed upstream behaviour (translation report chapter 7).

\begin{simctable}{small}{4}{>{\raggedright\arraybackslash}p{\dimexpr 0.372\linewidth-2\tabcolsep\relax}>{\raggedright\arraybackslash}p{\dimexpr 0.105\linewidth-2\tabcolsep\relax}>{\raggedright\arraybackslash}p{\dimexpr 0.257\linewidth-2\tabcolsep\relax}>{\raggedright\arraybackslash}p{\dimexpr 0.178\linewidth-2\tabcolsep\relax}}
%
\toprule\rowcolor{simchead}\simcth{id} & \simcth{Category} & \simcth{Fortran location (294f662)} & \simcth{upstream fix} \\\midrule\endfirsthead
\multicolumn{4}{@{}l}{\simccontinued}\\\toprule\rowcolor{simchead}\simcth{id} & \simcth{Category} & \simcth{Fortran location (294f662)} & \simcth{upstream fix} \\\midrule\endhead
\midrule\multicolumn{4}{r@{}}{\simccontinues}\\\endfoot
\bottomrule\endlastfoot
D-006 kaon decay, fixed upstream in 1fae2d7 (D-081) & A & shared/\allowbreak{}transp.\allowbreak{}f:\allowbreak{}158 & commit 1fae2d7 \\
D-038 F1 Δ++ STOP, fixed upstream in 20d46fc (D-081) & A & dbase.\allowbreak{}f:\allowbreak{}823-826 & commit 20d46fc \\
D-051 F3 pfermi NaN, fixed upstream in e4dd175 (D-081; since U.1 also in fortran\_\allowbreak{}compat) & A & results\_\allowbreak{}write.\allowbreak{}f:\allowbreak{}208-209 & commit e4dd175 \\
\end{simctable}
